\onecolumn
\lstset{%
  language=[LaTeX]TeX,
  backgroundcolor=\color{gray!25},
  basicstyle=\ttfamily,
  breaklines=true,
  breakindent=0pt,
  columns=fullflexible
}


\begin{lstlisting}[caption={Demonstration used for math aqua}, label=math aqua,,captionpos=b]
[Example 1]
[Instruction and Question] 
A person is traveling at 20 km/hr and reached his destiny in 2.5 hr then find the distance?

[Options]
(A) 53 km
(B) 55 km
(C) 52 km
(D) 60 km
(E) 50 km

[Answer] 
(E)

[Rationale]
The distance that the person traveled would have been 20 km/hr * 2.5 hrs = 50 km. So the answer is (E). 



[Example 2]
[Instruction and Question] 
John found that the average of 15 numbers is 40. If 10 is added to each number then the mean of the numbers is?

[Options]
(A) 50
(B) 45
(C) 65
(D) 78
(E) 64

[Answer] 
(A)

[Rationale]
If 10 is added to each number, then the mean of the numbers also increases by 10. So the new mean would be 50.  So the answer is (A).



[Example 3]
[Instruction and Question] 
If a / b = 3/4 and 8a + 5b = 22,then find the value of a.

[Options]
(A) 1/2
(B) 3/2
(C) 5/2
(D) 4/2
(E) 7/2

[Answer] 
(B)

[Rationale]
If a / b = 3/4, then b = 4a / 3. So 8a + 5(4a / 3) = 22. This simplifies to 8a + 20a / 3 = 22, which means 44a / 3 = 22. So a is equal to 3/2.  So the answer is (B).



[Example 4]
[Instruction and Question] 
How many keystrokes are needed to type the numbers from 1 to 500?

[Options]
(A) 1156
(B) 1392
(C) 1480
(D) 1562
(E) 1788

[Answer] 
(B).

[Rationale]
There are 9 one-digit numbers from 1 to 9. There are 90 two-digit numbers from 10 to 99. There are 401 three-digit numbers from 100 to 500. 9 + 90(2) + 401(3) = 1392.  So the answer is (B).



[Example 5]
[Instruction and Question]

    \end{lstlisting}
    




\begin{lstlisting}[caption={Demonstration used for math prompt}, label=math prompt,,captionpos=b]
[Example 1]
[Instruction and Question] 
There are 15 trees in the grove. Grove workers will plant trees in the grove today. After they are done, there will be 21 trees. How many trees did the grove workers plant today?

[Answer]
6

[Rationale]
We start with 15 trees. Later we have 21 trees. The difference must be the number of trees they planted. So, they must have planted 21 - 15 = 6 trees.



[Example 2]
[Instruction and Question] 
If there are 3 cars in the parking lot and 2 more cars arrive, how many cars are in the parking lot?

[Answer] 
5

[Rationale]
There are 3 cars in the parking lot already. 2 more arrive. Now there are 3 + 2 = 5 cars. 



[Example 3]
[Instruction and Question] 
Leah had 32 chocolates and her sister had 42. If they ate 35, how many pieces do they have left in total?

[Answer] 
39

[Rationale]
Leah had 32 chocolates and Leah's sister had 42. That means there were originally 32 + 42 = 74 chocolates. 35 have been eaten. So in total they still have 74 - 35 = 39 chocolates. 



[Example 4]
[Instruction and Question] 
Jason had 20 lollipops. He gave Denny some lollipops. Now Jason has 12 lollipops. How many lollipops did Jason give to Denny?

[Answer] 
8

[Rationale]
Jason had 20 lollipops. Since he only has 12 now, he must have given the rest to Denny. The number of lollipops he has given to Denny must have been 20 - 12 = 8 lollipops. 



[Example 5]
[Instruction and Question] 
Shawn has five toys. For Christmas, he got two toys each from his mom and dad. How many toys does he have now?

[Answer] 
9

[Rationale]
He has 5 toys. He got 2 from mom, so after that he has 5 + 2 = 7 toys. Then he got 2 more from dad, so in total he has 7 + 2 = 9 toys. 



[Example 6]
[Instruction and Question] 
There were nine computers in the server room. Five more computers were installed each day, from monday to thursday. How many computers are now in the server room?

[Answer] 
29

[Rationale]
There are 4 days from monday to thursday. 5 computers were added each day. That means in total 4 * 5 = 20 computers were added. There were 9 computers in the beginning, so now there are 9 + 20 = 29 computers. 



[Example 7]
[Instruction and Question] 
Michael had 58 golf balls. On tuesday, he lost 23 golf balls. On wednesday, he lost 2 more. How many golf balls did he have at the end of wednesday?

[Answer] 
33

[Rationale]
Michael initially had 58 balls. He lost 23 on Tuesday, so after that he has 58 - 23 = 35 balls. On Wednesday he lost 2 more so now he has 35 - 2 = 33 balls. 



[Example 8]
[Instruction and Question] 
Olivia has $23. She bought five bagels for $3 each. How much money does she have left?


[Answer] 
8

[Rationale]
She bought 5 bagels for $3 each. This means she spent 5 * $3 = $15 on the bagels. She had $23 in beginning, so now she has $23 - $15 = $8. 



[Example 9]
[Instruction and Question]

    \end{lstlisting}
    




\begin{lstlisting}[caption={Demonstration used for misc cola}, label=misc cola,,captionpos=b]
[Example 1]
[Instruction and Question]
Is the following sentence linguistically acceptable?\nFred must both have been singing songs and have been drinking beer.\nOPTIONS:\n- unacceptable\n- acceptable\n\n

[Answer]
acceptable

[Rationale]
The given sentence follows the subject-verb-object structure. The subject is "puppy," the verb is "was kissed," and the object is "by the policeman." Therefore, the sentence is linguistically acceptable.



[Example 2]
[Instruction and Question]
Problem: Produce a sentence that would be considered grammatically acceptable (OPTIONS:\n- unacceptable\n- acceptable)\n\nAnswer: 

[Answer]
The hanging gardens are a sight to behold.

[Rationale]
Here's an example of a grammatically correct sentence:\nThe hanging gardens are a sight to behold.



[Example 3]
[Instruction and Question]
Sentence: "Amanda drove at the package."\nWould a linguist rate this sentence to be acceptable linguistically? OPTIONS:\n- unacceptable\n- acceptable\nA: 

[Answer]
unacceptable

[Rationale]
From the sentence "Amanda drove at the package", it is not clear what is meant by "drove at." The phrasal verb "drove at" can mean "to move towards with the intention of hitting," which does not make sense in the context of a package. It is possible that the intended meaning was "Amanda drove to the package" or "Amanda drove towards the package," which would make more sense.



[Example 4]
[Instruction and Question]
Sentence: Ann may spends her vacation in Italy.\nOPTIONS:\n- unacceptable\n- acceptable

[Answer]
unacceptable

[Rationale]
The given sentence includs a subject-verb agreement error. The subject of the sentence is "Ann," which is singular, and the verb "may spends" is in the third person singular form. However, "may" is a modal verb, and when used with another verb, the second verb should be in the base form (i.e., the infinitive form without "to").\nTherefore, the correct verb form should be "spend" instead of "spends," and the correct sentence would be "Ann may spend her vacation in Italy.



[Example 5]
[Instruction and Question]
Instruction: Produce a brief English sentence that would be considered grammatically as category: unacceptable\nAll categories: OPTIONS:\n- unacceptable\n- acceptable\nA brief sentence: 

[Answer]
That monkeys is eating the banana.

[Rationale]
Here's an example of an unacceptable English sentence: That monkeys is eating the banana.



[Example 6]
[Instruction and Question]

    \end{lstlisting}
    




\begin{lstlisting}[caption={Demonstration used for misc math dataset}, label=misc math dataset,,captionpos=b]
[Example 1]
[Instruction and Question]
question: Teacher asked me this: Solve 3*j + 16 = -5*j for j.\n++++++++++\nanswer:

[Answer]
-2

[Rationale]
The given equation is 3j + 16 = -5j. To solve for j, we need to isolate j on one side of the equation.

First, we can simplify the equation by adding 5j to both sides:

3j + 16 + 5j = 0

Combining like terms, we get:

8j + 16 = 0

Subtracting 16 from both sides:

8j = -16

Dividing both sides by 8:

j = -2

Therefore, the solution for j is -2.



[Example 2]
[Instruction and Question]
Solve 1875954*s - 817 = 1875935*s for s.\nSolve this plz.\nA:

[Answer]
43

[Rationale]
We need to isolate the variable s on one side of the equation.

[STEP 1] 1875954*s - 1875935*s = 817
[STEP 2] 19*s = 817
[STEP3] s = 817/19
[STEP4] s = 43



[Example 3]
[Instruction and Question]
QUESTION: Math problem: Solve 2767686 = 158*f + 2754414 for f.\nANS:

[Answer]
84

[Rationale]
[STEP 1] 2767686 - 2754414 = 158*f
[STEP 2] 13272 = 158*f
[STEP 3] 13272/158 = f
[STEP 4] f = 84



[Example 4]
[Instruction and Question]
Q: What is the solution?\nSolve -46*y + 24 = -34*y for y.\n\nA:

[Answer]
2

[Rationale]
We need to isolate the variable y on one side of the equation.

[STEP 1] 46*y - 34*y = 24
[STEP 2] 12*y = 24
[STEP3] y = 24 / 12
[STEP4] y = 2



[Example 5]
[Instruction and Question]
Write down the solution for this math problem: Solve -21181 + 8447 = 497*p + 9401 + 3709 for p.\nanswer:

[Answer]
-52

[Rationale]
-21181 + 8447 = 497p + 9401 + 3709

Simplifying the left-hand side of the equation:

-12734 = 497p + 13110

Subtracting 13110 from both sides:

-25844 = 497p

Dividing both sides by 497:

p = -52

Therefore, the solution for the given math problem is p = -52.



[Example 6]
[Instruction and Question]
Q: Solve 493 = 18*r + 403 for r.\n\nA:

[Answer]
5

[Rationale]
To solve for r, we need to isolate the variable on one side of the equation.

Starting with:

493 = 18*r + 403

Subtract 403 from both sides:

90 = 18*r

Divide both sides by 18:

r = 5

Therefore, the solution for r is 5.



[Example 7]
[Instruction and Question]

    \end{lstlisting}
    




\begin{lstlisting}[caption={Demonstration used for nli prompt}, label=nli prompt,,captionpos=b]
[Example 1]
[Instruction and Question]
Problem: Context: Anorexia in males accounted for approximately six percent of cases seen in an eating disorder clinic.\nAnorexia becomes common in males.\nOPTIONS:\n- yes\n- no\nA: 

[Answer]
no

[Rationale]
The information given only pertains to cases seen in this particular eating disorder clinic and may not be representative of the entire population. To determine whether anorexia is becoming common in males, we would need to look at data from a larger sample size or population, as well as data from multiple sources to verify the trend. So the answer is “no”.



[Example 2]
[Instruction and Question]
Generate a context and a hypothesis.\n\nAnswer: 

[Answer]
He also liked swimming and cycling. He said that he wrote stories, though he had to admit that he had never got further than the first two pages. Willie meanwhile not only remained silent during these conversations but picked his berries slowly so that they might forget that he was there but he reckoned without Zach.\n\nHypothesis: Willie was there

[Rationale]
Here's a context and an example for natural language inference task:\nHe also liked swimming and cycling. He said that he wrote stories, though he had to admit that he had never got further than the first two pages. Willie meanwhile not only remained silent during these conversations but picked his berries slowly so that they might forget that he was there but he reckoned without Zach.\n\nHypothesis: Willie was there



[Example 3]
[Instruction and Question]
Q: Brock was a popular English writer and director.\nJeremy Brock MBE (born 1959) is an English writer and director whose works include the screenplays "Mrs Brown", "Driving Lessons", "The Last King of Scotland", "Charlotte Gray", and "The Eagle". Brock has also written two plays for the Hampstead downstairs theatre.\nOPTIONS:\n- Yes\n- It's impossible to say\n- No\nA: 

[Answer]
It's impossible to say

[Rationale]
While the information provided in the options about Jeremy Brock suggests that he is indeed an English writer and director with a number of successful works under his belt, there is not enough information to determine whether he is actually "popular" or not. The term "popular" is subjective and can vary depending on different factors, such as time period, genre, and audience. So the answer is "It's impossible to say".



[Example 4]
[Instruction and Question]
Biff Burger (Best In Fast Food) was a burger fast food franchise that originated in Clearwater, Florida. Early success in the late 1950s and early 1960s led the franchise to spread to many parts of the Southeastern United States. However, financial woes in the 1970s sent the chain on a decline and most of its restaurants out of business by the 1980s. Today, two Biff Burger restaurants remain.\nBiff Burger came about in Clearwater Florida. OPTIONS:\n- Yes\n- It's impossible to say\n- No\nA: 

[Answer]
Yes

[Rationale]
To validate the claim "Biff Burger came about in Clearwater Florida", we should find a supporting evidence within the context. We could see that "Biff Burger (Best In Fast Food) was a burger fast food franchise that originated in Clearwater, Florida." supports the claim we are trying to validate. So the answer is "Yes".



[Example 5]
[Instruction and Question]
input hypothesis: In the decade before the 1950s, the United States FBI began to maintain a public list of the people it regarded as the Ten Most Wanted Fugitives\nContext: In the 1950s, the United States FBI began to maintain a public list of the people it regarded as the Ten Most Wanted Fugitives. Following is a brief review of FBI people and events that place the 1950s decade in context, and then an historical list of individual fugitives whose names first appeared on the 10 Most Wanted list during the decade of the 1950s, under FBI Director J. Edgar Hoover.\nOPTIONS:\n- Yes\n- It's impossible to say\n- No\ntrue or false: 

[Answer]
No

[Rationale]
The context provided contradicts this hypothesis by stating that the FBI began to maintain this public list in the 1950s, under the direction of J. Edgar Hoover. The context goes on to list individuals who were first placed on the 10 Most Wanted list during the 1950s, further supporting the notion that the list was not in existence prior to that decade. So the answer is "No".



[Example 6]
[Instruction and Question]
Q: If The plaintiffs claim Penrose's design is distinguished by its aperiodicity (the pattern almost but never quite repeats itself) and its five-fold symmetry (a trait that at the time was thought not to exist in nature but has since been identified in certain crystal formations)., can we say that "The design has no identified traits that would distinguish it."?\nOPTIONS:\n- yes\n- it is not possible to tell\n- no\n\nA: 

[Answer]
no

[Rationale]
The plaintiffs are stating that the design has unique characteristics that set it apart, which directly contradicts the statement that the design has no identified traits that would distinguish it. Therefore, the answer is "no".



[Example 7]
[Instruction and Question]
Does "The first docking of Soviet spacecraft was finally realised in January 1969 by the Soyuz 4 and Soyuz 5 missions." contain the correct answer to "When was the first successful docking of a two man space crew?"\nOPTIONS:\n- yes\n- no\n\nA:

[Answer]
yes

[Rationale]
The passage talks about the first successful docking of a two-man space crew, as both Soyuz 4 and Soyuz 5 had two cosmonauts on board at the time of their docking. Therefore, the statement contains the correct answer to the question. So, the answer is "yes".



[Example 8]
[Instruction and Question]
input question: Is it true that "I shall have to guard The foxes." if "The foxes are getting in at night and attacking the chickens. I shall have to guard them." is true? OPTIONS:\n- no\n- yes\noutput answer: 

[Answer]
no

[Rationale]
The statement "I shall have to guard The foxes" does not accurately reflect the situation described in the second sentence "I shall have to guard them". In this context, "them" refers to the foxes, so the correct statement would be "I shall have to guard them", not "I shall have to guard the foxes". So, the answer is "no".



[Example 9]
[Instruction and Question]
Question:\nPremise: A man in a blue shirt and orange shorts looks on at a small child wearing a yellow shirt and blue shorts with trees comprising the backdrop.\n\nHypothesis: There are a lot of trees behind the small child.\nIs the hypothesis true?\nOPTIONS:\n- yes\n- it is not possible to tell\n- no\n\nAnswer:\n

[Answer]
it is not possible to tell

[Rationale]
While it is mentioned that trees make up the backdrop, it is not specified whether there are a lot of trees or not. Additionally, the focus of the premise is on the man and child and not on the trees in the background. Hence, it is not possible to determine whether the hypothesis is true or not.



[Example 10]
[Instruction and Question]

    \end{lstlisting}
    




\begin{lstlisting}[caption={Demonstration used for rc prompt}, label=rc prompt,,captionpos=b]
[Example 1]
[Instruction and Question]
Extract the answer to the question from the following context.
Question: Architects have to not only be talented designers, but at the other end of the spectrum, have a talent for the more mundane what?
Context: Architecture has to do with planning and designing form, space and ambience to reflect functional, technical, social, environmental and aesthetic considerations. It requires the creative manipulation and coordination of materials and technology, and of light and shadow. Often, conflicting requirements must be resolved. The practice of Architecture also encompasses the pragmatic aspects of realizing buildings and structures, including scheduling, cost estimation and construction administration. Documentation produced by architects, typically drawings, plans and technical specifications, defines the structure and/or behavior of a building or other kind of system that is to be or has been constructed.

[Options]
Not Given

[Answer]
pragmatic aspects of realizing buildings and structures

[Rationale]
The context states architects must have a talent for in addition to design. This includes activities such as scheduling, cost estimation and construction administration which are essential to creating a physical structure. 

Which is why the answer extracted from the context is "pragmatic aspects of realizing buildings and structures".



[Example 2]
[Instruction and Question]
Please answer correctly the following question related to the paragraph below. 
Which tribe was less affected by the disruption of ocean food webs?

Two related native tribes live thousands of miles away from each other. Ana tribe lives near the South Pole in Antarctica, while Momu tribe lives in Northern Alaska near the North Pole.

Hint: Unfortunately, the layer of good ozone is being destroyed by air pollution. The chief culprits are chlorine and bromine gases. They are released in aerosol sprays, coolants, and other products. Loss of ozone has created an ozone hole over Antarctica. Ozone depletion results in higher levels of UV radiation reaching Earth. In humans, this increases skin cancers and eye cataracts. It also disturbs the nitrogen cycle, kills plankton, and disrupts ocean food webs. The total loss of the ozone layer would be devastating to most life. It’s rate of loss has slowed with restrictions on pollutants, but it is still at risk.

[Options]
Not Given

[Answer]
Momu

[Rationale]
The paragraph mentions the negative impacts of the loss of the ozone layer, specifically the disruption of ocean food webs. The Ana tribe lives near the South Pole in Antarctica, where the ozone hole is located, while the Momu tribe lives near the North Pole.

Thus, based on the context of the ozone depletion affecting the ocean food webs, it can be inferred that the Momu tribe was less affected by the disruption of ocean food webs. 

Therefore, the answer is Momu.



[Example 3]
[Instruction and Question]
Given the below context:

Feeling unsettled, Etty left Rome for Venice, intending to remain there for 10 days and then return to England. Evans preferred to remain in Rome, so Etty travelled alone, pausing briefly in Florence and in Ferrara (where he stopped to kiss the armchair of Ludovico Ariosto). The painter Charles Lock Eastlake, then resident in Rome, had provided Etty with a letter of introduction to Harry D'Orville, British Vice consul in Venice; D'Orville was so impressed with Etty that he arranged for him to stay in his own house, rather than in lodgings. Etty had long considered Venice his spiritual home and "The hope and idol of my professional life", and had often wondered why, given its artistic importance, so few English travellers visited the city. He was not disappointed. Throughout the remainder of his life, he looked back on his visit to Venice with great fondness, writing shortly before his death that "Venezia, cara Venezia! thy pictured glories haunt my fancy now!"
Although Etty had only intended to stay for 10 days, he was so taken with Venice that he remained for over seven months. He fell into a routine of copying paintings in Venetian collections by day, and attending the life class of the Venetian Academy of Fine Arts by night, producing around 50 oil paintings in total as well as numerous pencil sketches. He was extremely impressed with the high quality of the Venetian Academy; the instructors in their turn were extremely impressed with the quality of Etty's work, in particular his flesh tones. He acquired the nickname of ""Il Diavolo"" owing to the high speed at which he was able to paint, and watching him at work became something of a spectacle in its own right; luminaries including Gioachino Rossini and Ladislaus Pyrker (then Patriarch of Venice) came to watch him paint. 
So devoted was Etty to his studies in Venice that he exhibited no original work in 1823, writing to his brother that "If one spent all the time in painting originals, one might as well, nay better, be at home". The members of the Venetian...

Guess a valid title for it!

[Options]
Not Given

[Answer]
William Etty

[Rationale]
The context is a description of the experiences and events surrounding the British painter William Etty during his visit to Venice. Throughout the text, the focus is on Etty and his actions, thoughts, and feelings, making it clear that he is the main subject of the paragraph. The text mentions his nickname ""Il Diavolo"" and the admiration he received from notable figures like Gioachino Rossini and Ladislaus Pyrker, further emphasizing his central role.

Given the subject matter and focus on Etty, a fitting title for the paragraph could be "William Etty".



[Example 4]
[Instruction and Question]
I want to test the ability of students to read a passage and answer questions about it. Could you please come up with a good question for the passage?
"Some high-speed black-and-white films, such as Ilford Delta 3200 and Kodak T-MAX P3200, are marketed with film speeds in excess of their true ISO speed as determined using the ISO testing method. For example, the Ilford product is actually an ISO 1000 film, according to its data sheet. The manufacturers do not indicate that the 3200 number is an ISO rating on their packaging. Kodak and Fuji also marketed E6 films designed for pushing (hence the "P" prefix), such as Ektachrome P800/1600 and Fujichrome P1600, both with a base speed of ISO 400."

[Options]
Not Given

[Answer]
What is notable about these films?

[Rationale]
The key point that the passage emphasizes is the difference between the marketed film speeds and the actual ISO speeds of these films.
The passage raises questions about the marketing practices in the photography industry and the reliability of the information provided to consumers.The manufacturers of these films do not indicate the actual ISO ratings on the packaging, which makes it notable. The passage also mentions that these films are marketed for pushing, which is another notable aspect of these films.
Thus, the best answer would be: What is notable about these films?"



[Example 5]
[Instruction and Question]
I am a movie director and I just received the following movie plot. 
Could you help me answer this question? If not, let me know by writing "Not answerable".

Plot title: Bananas
Movie plot: Fielding Mellish (Woody Allen) is the protagonist, but he does not appear until after the opening credits. The cold open, which featured the assassination of the president of the fictional ""banana republic"" of San Marcos and a coup d'état that brings Gen. Emilio Molina Vargas (Carlos Montalban) to power, sets up the situation in which Mellish would become embroiled. The scene was in the form of a championship boxing telecast on Wide World of Sports, with Don Dunphy as the host and Howard Cosell as the commentator.[4] Mellish is a neurotic blue collar man who tries to impress social activist Nancy (Louise Lasser) by trying to get in touch with the revolution in San Marcos. He visits the republic and attempts to show his concern for the native people. However, he is nearly killed by the local caudillo and then saved by the revolutionaries, putting him in their debt. Mellish then learns, clumsily, how to be a revolutionary. When the revolution is successful, the Castro-style leader goes mad, forcing the rebels to place Mellish as their President. When traveling back to the U.S. to obtain financial aid, he reunites with his activist ex-girlfriend and is exposed. In a classic courtroom scene, Mellish tries to defend himself from a series of incriminating witnesses, including a reigning Miss America and a middle-aged African-American woman claiming to be J. Edgar Hoover in disguise. One of the witnesses does provide testimony favorable to Mellish, but the court clerk, when asked to read back this testimony, replies with an entirely different, wholly unfavorable rendition. 
Mellish is eventually sentenced to prison, but his sentence is suspended on the condition that he does not move into the judge's neighborhood. Nancy then agrees to marry him. The film ends with the between-the-covers consummation of their marriage, an event that was over much more quickly than Nancy had anticipated. Like the opening scene, it was accompanied by Cosell providing commentary.

My question: What is Feilding Mellish's occupation?"

[Options]
Not Given

[Answer]
Blue collar worker who becomes president of San Marcos

[Rationale]
In the movie plot, Feilding Mellish is described as a "neurotic blue collar man". The term "blue collar" refers to a type of job that is typically manual labor or involves working with machinery. Blue collar workers often perform physical labor, as opposed to white collar workers who perform administrative, professional, or managerial work. 
So, based on this description, we can conclude that Feilding Mellish's occupation at the start of the movie is a blue-collar worker. Later in the movie, Feilding Mellish becomes involved in a revolution in the fictional banana republic of San Marcos and is eventually placed as its president after the previous leader goes mad.
Therefore, the answer to the question is: "Blue collar worker who becomes President of San Marcos".



[Example 6]
[Instruction and Question]
This article: 
Shackleton's February 1907 announcement that he intended to base his expedition at the old Discovery headquarters was noted by Scott, whose own future Antarctic plans were at that stage unannounced. In a letter to Shackleton, Scott claimed priority rights to McMurdo Sound. "I feel I have a sort of right to my own field of work," he wrote, adding: "anyone who has had to do with exploration will regard this region primarily as mine".He concluded by reminding Shackleton of his duty of loyalty towards his former commander.
Shackleton's initial reply was accommodating: "I would like to fall in with your views as far as possible without creating a position that would be untenable to myself". Edward Wilson, asked by Shackleton to mediate, took an even tougher line than Scott. 
"I think you should retire from McMurdo Sound", he wrote, advising Shackleton not to make any plans to work from anywhere in the entire Ross Sea quarter until Scott decided 
""what limits he puts on his own rights"". 
To this Shackleton replied: "There is no doubt in my mind that his rights end at the base he asked for [...] I consider I have reached my limit and I go no further".The matter was unresolved when Scott returned from sea duty in May 1907. Scott pressed for a line of demarcation at 170° W—everything to the west of that line, including Ross Island, McMurdo Sound, and Victoria Land, would be Scott's preserve. Shackleton, with other concerns pressing on him, felt obliged to concede. 
On 17 May he signed a declaration stating that ""I am leaving the McMurdo base to you"", and that he would seek to land further east, either at the Barrier Inlet visited briefly during the Discovery Expedition, or at King Edward VII Land. He would not touch the coast of Victoria Land at all. It was a capitulation to Scott and Wilson, and meant forfeiting the expedition's aim of reaching the South Magnetic Pole which was located within Victoria Land. 
Polar historian Beau Riffenburgh believes this was "a promise that should never ethically have been demanded and... 

contains an answer for the question: 
"What are the specific names of the two men between whom the dispute soured relations?", what is it ?

[Options]
Not Given

[Answer]
Shackleton

[Rationale]
From the given context, the answer obviously seems to be: Shackleton



[Example 7]
[Instruction and Question]

    \end{lstlisting}
    




\begin{lstlisting}[caption={Demonstration used for rc drop}, label=rc drop,,captionpos=b]

    \end{lstlisting}
    




\begin{lstlisting}[caption={Demonstration used for cbqa prompt}, label=cbqa prompt,,captionpos=b]
[Example 1]
[Instruction and Question]
Here's a complex question that requires someone to reason about the input, can you answer it?
In what country were the second Summer Olympics at which Hanna Titimets represented Ukraine celebrated?

[Options]
Not Given

[Answer]
Brazil

[Rationale]
Hanna Titimets represented Ukraine in the 2016 Summer Olympics which were held in Rio de Janeiro, Brazil. This was the second Summer Olympics where Ukraine was represented and the only one held in Brazil.



[Example 2]
[Instruction and Question]
Question: what is modified agi?
Would "It includes wages, interest, dividends, business income, rental income, and all other types of income." be a reasonable answer?

[Options]
- Yes
- No

[Answer]
No

[Rationale]
"Modified AGI" refers to Adjusted Gross Income, which is a term used in taxation. It is calculated as a person's total income, minus certain deductions that are allowed by the tax code. 
The answer provided, "It includes wages, interest, dividends, business income, rental income, and all other types of income" describes the types of income that are included in a person's total income, but it does not describe Modified AGI specifically. 
In conclusion, based on this information, the answer provided is not a reasonable answer to the question of what Modified AGI is, so the answer is "No".



[Example 3]
[Instruction and Question]
Answer this question: what is renaissance english?

[Options]
Not Given

[Answer]
The English Renaissance was a cultural and artistic movement in England dating from the late 15th and early 16th centuries to the early 17th century.

[Rationale]
Here's an direct answer: The English Renaissance was a cultural and artistic movement in England dating from the late 15th and early 16th centuries to the early 17th century.



[Example 4]
[Instruction and Question]
Combine facts and answer this: The 2009 Tostitos Fiesta Bowl game was a post-season college football bowl game between the Ohio State Buckeyes and the Texas Longhorns on Monday, January 5, 2009, the Big 12 champion University of Oklahoma Sooners were participating in the 2009 BCS National Championship Game, was an American football game played at Dolphin Stadium in Miami Gardens, Florida, on which date?

[Options]
Not Given

[Answer]
January 8, 2009

[Rationale]
The 2009 FedEx BCS National Championship Game was an American football game played at Dolphin Stadium in Miami Gardens, Florida, on January 8, 2009.



[Example 5]
[Instruction and Question]
FINAL EXAM 
Question 1. Matt's Bar in Minneapolis claims to have created what cheeseburger that has the cheese inside the meat patty?

[Options]
Not Given

[Answer]
Jucy Lucy

[Rationale]
The answer is based on the claim made by Matt's Bar in Minneapolis that they have created the cheeseburger with the cheese inside the meat patty, which is known as the Jucy Lucy. 
It is called Jucy Lucy because the cheese is inside the meat patty causing it to be juicy. 

Therefore, the answer should be "Jucy Lucy".



[Example 6]
[Instruction and Question]
I am verifying the answers generated by an automatic system to the following question: what the names of destiny's child
Suggested answer: Destiny's Child was an American R&B girl group whose final, and perhaps most recognizable, line-up comprised Beyonc\u00e9 Knowles , Kelly Rowland and Michelle Williams .

Should I validate this answer?

[Options]
- Yes
- No

[Answer]
Yes

[Rationale]
The statement that "Destiny's Child was an American R&B girl group whose final and perhaps most recognizable line-up comprised of Beyoncé Knowles, Kelly Rowland, and Michelle Williams" is a factual representation of the group. This information can be verified through various credible sources such as music history books, biographies, and articles dedicated to the group and its members. Additionally, the group's recognition and popularity further support the validity of the suggested answer.
The suggested answer provided by the automatic system is accurate and should be validated. 
Therefore, the answer is Yes.



[Example 7]
[Instruction and Question]

    \end{lstlisting}




\begin{lstlisting}[caption={Demonstration used for commonsense prompt}, label=commonsense prompt,,captionpos=b]
[Example 1]
[Instruction and Question]
Skylar had stopped breathing but Lee held Skylar's arms anyway, in vain. Given the context: Lee want to do what next?
Possible answers: beg the doctors to try again, get her up and walking home, bring Skylar back

[Options]
- beg the doctors to try again
- get her up and walking home
- bring Skylar back

[Answer]
beg the doctors to try again

[Rationale]
The context of the situation is that Skylar has stopped breathing and Lee is holding Skylar's arms. This suggests that Skylar has received medical attention and is no longer alive.
Based on this, Lee would want to beg the doctors to try again, as this would be the best way to potentially save Skylar's life. 
The options of getting Skylar up and walking home or bringing Skylar back do not align with the context, as they would not address the immediate medical emergency at hand.

So the answer is to beg the doctors to try again.



[Example 2]
[Instruction and Question]
Do you think the right answer to the question "what can run on alcoholic fermentation of the glucose in corn or other plants?" is "bicycles"?

[Options]
- No
- Yes

[Answer]
No

[Rationale]
Alcoholic fermentation is a process that occurs in living organisms, most commonly yeast, to convert glucose into energy. This process does not produce enough energy to power something like a bicycle.

So the answer is No.



[Example 3]
[Instruction and Question]
Eventually , the cyst can get so large that part of the wall weakens and fluid begins to leak out ; that is , the cyst ruptures . It is similar to popping a water balloon . When this happens , the fluid may irritate the pelvic lining and cause pain , often severe . Most of the time the pain begins to gradually subside over the course of the next few days .
What type of doctor might the person have had to see ?

- A paediatrician  - A heart surgeon - a gynecologist - None of the above choices

[Options]
- A paediatrician 
- A heart surgeon 
- a gynecologist 
- None of the above choices

[Answer]
a gynecologist

[Rationale]
The context mentions a cyst in the pelvic region and the symptoms of the cyst rupturing, including pain and fluid leakage. This suggests that the issue is related to the reproductive or urinary system. Given that the issue is related to the reproductive or urinary system, it is likely that the person would have had to see a specialist in these areas. 
Out of the given options, "a gynecologist" is the most likely specialist that the person would have had to see, as gynecologists specialize in female reproductive health, including issues related to the uterus, ovaries, and other reproductive organs. 

In conclusion, based on the given context, the most likely answer is a gynecologist.



[Example 4]
[Instruction and Question]
Process:
- Get some seeds
- Pick a spot to plant them
- Dig a hole in the dirt
- Put the seed in the hole
- Pour some water on the seed and hole
- Cover up the hole
- Press down on it
- Spray some plant food on it.

Question:
suppose if no water is available happens, how will it affect vegetables?

- A: more
- B: less

[Options]
- A
- B

[Answer]
B

[Rationale]
The process of planting seeds requires water, which is essential for the growth of plants. Water helps transport nutrients, regulate temperature, and absorb minerals from the soil. If no water is available, the seed will not germinate and the plant will not grow properly, leading to fewer vegetables. Hence, the absence of water will result in "less" vegetables, making the answer "B".



[Example 5]
[Instruction and Question]
Read the following conversation and answer the question.
M: I'm worried about my sister, Lucy.
W: Why is that?
M: She thinks she's too fat.
W: Is she?
M: No, but she keeps skipping meals. Then, she only eats chips and drinks cola.
W: I used to do that. It's called binge. It was no fun!
M: Why did you stop doing it?
W: Well, my doctor told me to eat when I'm hungry. She said, "Eat till you're full or you'll eat too much later." She said a lot of girls ruin their health this way.
M: Did she say what to eat?
W: She said, "Eat fruit, vegetables, meats and grains. Have regular meals and snacks. Get exercise, too."

Question: What's the relationship between the speakers? 

- Doctor and patient.
- Friends.
- Brother and sister.

[Options]
- Doctor and patient
- Friends
- Brother and sister

[Answer]
Friends

[Rationale]
The conversation between two people, a man(M) and a woman(W), is discussing the man's sister who has a negative view of her own body and is skipping meals. The woman shares that she used to engage in similar behavior, and explains that her doctor advised her to eat a balanced diet and exercise.

Among the three options,
- First, since the woman is sharing her experience of what the doctor told her, it is unlikely she would be a medical professional. So the answer couldn't be doctor and patient.
- Second, given that the man is talking about his sister, the woman is unlikely to be the man's sister.
- Lastly, because the conversation is informal and the speakers are discussing a personal matter it is likely that the relationship between the speakers is that of friends.

Therefore, the answer is Friends.



[Example 6]
[Instruction and Question]
Question: The best place for the troublemaker to sneak a smoke while he was in where, was in the bathroom?

Choices: 
- playground
- school
- at hotel
- theater
- convenience store

The rationale to choose ""school"" as the answer is that:

[Options]
Not Given

[Answer]
smokers traditionally use school bathrooms

[Rationale]
Here's the rationale: "smokers traditionally use school bathrooms".



[Example 7]
[Instruction and Question]
So I was looking at quotes for our tshirts for ice hockey , and I came across one that sums up how I ' m feeling right now . When you talk about an injury and the kind of depression you go through , it 's not just because you 're out of shape and you can't go out and play . You 're missing a part of you .

Why is the narrator so down in the dumps ?
Pick the best answer from the following options:
A. They are unable to play ice curling .
B. They are hurt .
C. They are unable to play ice skating .
D. They are unable to play in the yard .

[Options]
- A
- B
- C
- D

[Answer]
B

[Rationale]
The quote in the question mentions the narrator's feelings about an injury and the depression that follows. The narrator says that it's not just about being out of shape and unable to play, but rather they feel like they are missing a part of themselves. This implies that the reason for their depression is linked to their physical inability to participate in a sport or activity they enjoy. Based on the context, it can be inferred that the narrator is upset because they are hurt and unable to play ice hockey.
So, the answer is B.



[Example 8]
[Instruction and Question]
Q: What does the ubiquitous generator use to generate currents?

Choices:

- inertia - water - gravity - magnetism

[Options]
- inertia
- water
- gravity
- magnetism

[Answer]
magnetism

[Rationale]
A generator is a device that converts mechanical energy into electrical energy. This is done by using a magnetic field to induce a current in the generator’s windings. This current is then sent to an external circuit, where it can be used to power electrical devices. The generator uses magnetism to generate currents, so the answer is magnetism.



[Example 9]
[Instruction and Question]

    \end{lstlisting}
    




\begin{lstlisting}[caption={Demonstration used for exqa drop}, label=exqa drop,,captionpos=b]
[Example 1]
[Instruction and Question] 
Read this article and answer this question The Ravens started their season at home against the Bengals.  They jumped out to a 3-0 lead in the first quarter with Justin Tucker's 46-yard field goal.  This was followed up by Ray Rice's 7-yard run to make the score 10-0.  The Bengals got on the board in the 2nd quarter with Mike Nugent's 34-yard field goal to shorten the lead to 10-3.  However, the Ravens were able to pull away as Joe Flacco found Anquan Boldin on a 34-yard touchdown pass to move ahead 17-3.  The Bengals responded coming within 7 when Benjarvus Green-Ellis ran for a 6-yard touchdown making the score 17-10 at halftime.  In the 3rd quarter, the Bengals were able to get within 4 points with Nugent kicking a 19-yard field goal shortening the Ravens' lead to 4 17-13.  However, the Ravens overpowered the Bengals scoring 27 unanswered points as Flacco found Dennis Pitta on a 10-yard pass to move ahead 24-13 followed up by Tucker scoring a 40-yard field goal moving them ahead 27-13.  This was followed by an Ed Reed interception that was returned 34 yards for a touchdown moving the team ahead 34-13.  In the 4th quarter, the Ravens scored again off of Rice's 1-yard run for a 41-13 lead and finished the game off of Tucker's 39-yard field goal to make the final score 44-13. Ed Reed returned an interception for a 34-yard touchdown, making Reed the all-time leader in career interception return yards with 1,497. The previous record of 1,483 yards was held by Rod Woodson.\nHow many yards did Rice rush for on his two touchdown runs?

[Answer] 
8

[Rationale]
To calculate how far Rice ran in his two touchdown runs, you need to check how far he ran in each run. First, Rice ran 7-yards when the score became 10-0. Second, Rice ran 1-yard when the score became 41-13. Therefore, we could conclude that he ran 7+1=8 yards in total.



[Example 2]
[Instruction and Question]
The Colts came into this game three games ahead of the Jaguars in the AFC South standings, and with a win in this game, the Colts can clinch a playoff spot and the AFC South Championship for the fourth straight year. On the Jaguars' first play from scrimmage, RB Fred Taylor ran up the middle for 76 yards down to the Colts' 18-yard line. On the following play rookie RB Maurice Jones-Drew rushed 18 yards for the first score of the game. The Colts led the Jaguars at one point in the game, 10-7, but after Jacksonville scored six times before the Colts scored again, there was no way to catch up. The Colts allowed 375 rushing yards in this game, the second-highest total since the NFL-AFL merger in 1970. Jacksonville RB Maurice Jones-Drew ran for 166 yards and RB Fred Taylor ran for 131 yards. Third-string RB Alvin Pearman also ran for 71 yards. To further emphasize how effective the Jacksonville running game was, Jaguars QB David Garrard was only 8 for 14 with 79 yards passing. While he only threw the ball 14 times, Colts QB Peyton Manning threw the ball 50 times, completing 25 of those passes for 313 passing yards. Neither quarterback threw a touchdown pass, but both of them threw one interception each. The Colts WR tandem of Marvin Harrison and Reggie Wayne did well in this game. Harrison had 8 catches for 110 yards receiving, and Wayne had 6 catches for 101 yards. Jaguars RB Maurice Jones-Drew did not just succeed on offense&#8212;he ran back an Adam Vinatieri kickoff 93 yards for a touchdown also. The Colts lost SS Antoine Bethea to a shoulder injury, and he would not return. The Colts moved to 10-3, losing first place in the AFC, while the Jaguars improved to 8-5.\n\nAsk a question about this article.

[Answer]
how many yards did Taylor run up?

[Rationale]
Here's a question for the given article:\n\nQuestion: how many yards did Taylor run up?



[Example 3]
[Instruction and Question] 
Macquarie University (MQ) world rankings includes it being number 240 on the QS rankings, number 251+ on Times (THE), number 151+ on ARWU, and  number 267 with US News. This contributes to Macquarie being the number 8 ranked Australian university overall in the world ranking systems. Macquarie University rankings within Australia include being placed at number 8 on the ERA scale (2012) and being a 4 1/2 Star AEN rated university. Macquarie also has a student survey satisfaction rating of 77.4% for business, 90.3% for health, 91.4% for arts, and 93.8% for science. Macquarie is ranked in the top 40 universities in the Asia-Pacific region and within Australias top 12 universities according to the Academic Ranking of World Universities, the U.S. News & World Report Rankings and the QS World University Rankings. Macquarie was the highest ranked university in Australia under the age of 50 and was ranked 18th in the world (prior to its golden jubilee in 2014), according to the QS World University Rankings.\n\nHow many points higher is the satisfaction rating of science than arts?

[Answer] 
2.4

[Rationale]
In order to answer how many points higher the satisfaction rating of science is compared to arts, we first need to identify each satisfaction rating and subtract them. (STEP 1) science = 93.8%, (STEP 2) art = 91.4% (STEP 3) difference = 93.8 - 91.4 = 2.4%



[Example 4]
[Instruction and Question] 
Context: Samuel Shute was an English military officer and royal governor of the provinces of Massachusetts and New Hampshire.  After serving in the Nine Years' War and the War of the Spanish Succession, he was appointed by King George I as governor of Massachusetts and New Hampshire in 1716.  His tenure was marked by virulent disagreements with the Massachusetts assembly on a variety of issues, and by poorly conducted diplomacy with respect to the Native American Wabanaki Confederacy of northern New England that led to Dummer's War . Although Shute was partly responsible for the breakdown in negotiations with the Wabanakis, he returned to England in early 1723 to procure resolutions to his ongoing disagreements with the Massachusetts assembly, leaving conduct of the war to Lieutenant Governor William Dummer.  His protests resulted in the issuance in 1725 of the Explanatory Charter, essentially confirming his position in the disputes with the assembly.  He did not return to New England, being replaced as governor in 1728 by William Burnet, and refused to be considered for reappointment after Burnet's sudden death in 1729. Thomas Hutchinson , in his history of Massachusetts, described Shute's tenure as governor as the most contentious since the Antinomian Controversy of the 1630s.\n\nQuestion: Who was governor first: William Dummer or William Burnet?\n\nAnswer:

[Answer] 
William Dummer

[Rationale]
To solve the problem "Who was governor first: William Dummer or William Burnet?", we need to compare the year when William Dummer and William Burnet were governors. For William Dummer, it is mentioned that he was governor when Shute returned to England in early 1723. For William Burnet, it is mentioned that he became the governor in 1728. Since 1723 is earlier than 1728, William Dummer became the governor earlier than William Burnet.



[Example 5]
[Instruction and Question]
Write an article that answers the following question: How many cities did the refugees of the Hamidian massacres originate?

[Answer]
According to available information, the about 1,000 refugees from the Hamidian massacres  mainly originated from Diyarbakir , Aintab and Kilis; only about 100 of them stayed. The next wave of Armenian refugees were the about 2,000 who fled the Adana massacre in 1909, most of whom returned to their ancestral homes in Adana within the same year. However, the largest wave of Armenian refugees - some of whom had come before and returned - were the nearly 9,000 who escaped the massive deportations, the horrific massacres and the Genocide perpetrated by the Ottomans and the Young Turks; about 1,300 of them decided to stay, while the others eventually made arrangements to settle in other countries. Those refugees came mainly from Adana and Seleucia , while there a significant number of them came from Sis, Marash, Tarsus, Caesarea, Hadjin and Aintab; smaller numbers came from other places, alphabetically: Adapazar, Adrianople , Afion-Karahisar, Alexandretta , Arapgir, Armash, Baghche, Bardizag, Balian Dagh, Biredjik, Bitlis, Brusa, Chemishgezek, Constantinople , D\u00f6rtyol, Edessa ,  Erzerum, Eskishehir, Everek, Ikonion , Jeyhan, Kesab, Kharpert, Kutahia, Malatia, Mersin, Misis, Musa Dagh , Nicomedia , Rhaedestos , Sasun, Sebastia , Shar, Sivri Hisar, Smyrna , Tokat , Trepizond, Van, Yerzinga, Yozgat and Zeitun.

[Rationale]
Here's an example of an article for the given question:\n\nArticle: According to available information, the about 1,000 refugees from the Hamidian massacres  mainly originated from Diyarbakir , Aintab and Kilis; only about 100 of them stayed. The next wave of Armenian refugees were the about 2,000 who fled the Adana massacre in 1909, most of whom returned to their ancestral homes in Adana within the same year. However, the largest wave of Armenian refugees - some of whom had come before and returned - were the nearly 9,000 who escaped the massive deportations, the horrific massacres and the Genocide perpetrated by the Ottomans and the Young Turks; about 1,300 of them decided to stay, while the others eventually made arrangements to settle in other countries. Those refugees came mainly from Adana and Seleucia , while there a significant number of them came from Sis, Marash, Tarsus, Caesarea, Hadjin and Aintab; smaller numbers came from other places, alphabetically: Adapazar, Adrianople , Afion-Karahisar, Alexandretta , Arapgir, Armash, Baghche, Bardizag, Balian Dagh, Biredjik, Bitlis, Brusa, Chemishgezek, Constantinople , D\u00f6rtyol, Edessa ,  Erzerum, Eskishehir, Everek, Ikonion , Jeyhan, Kesab, Kharpert, Kutahia, Malatia, Mersin, Misis, Musa Dagh , Nicomedia , Rhaedestos , Sasun, Sebastia , Shar, Sivri Hisar, Smyrna , Tokat , Trepizond, Van, Yerzinga, Yozgat and Zeitun.



[Example 6]
[Instruction and Question]
In March 1768, Ming Rui began his retreat, pursued by a Burmese army of 10,000 men and 2000 cavalry. The Burmese then tried to encircle the Chinese by splitting the army into two. Maha Thiha Thura had now assumed the overall command, replacing Maha Sithu. The smaller army, led by Maha Sithu, continued to pursue Ming Rui while the larger army led by Maha Thiha Thura advanced through the mountainous route to emerge directly behind the Chinese. Through careful maneuvering, the Burmese managed to achieve complete encirclement of the Chinese at modern-day Pyinoolwin , about 50 miles northeast of Ava. Over the course of three days of bloody fighting, the Bannerman army was completely annihilated. The slaughter was such that the Burmese could hardly grip their swords as the hilts were slippery with enemy blood. Of the original 30,000 men of the main army, only 2500 remained alive and were captured. The rest had been killed either on the battlefield, through disease or through execution after their surrender. Ming Rui himself was severely wounded in battle. Only a small group managed to break through and escaped the carnage. Ming Rui himself could have escaped with that group. Instead, he cut off his queue and sent it to the emperor as a token of his loyalty by those who were escaping. He then hanged himself on a tree. In the end, only a few dozen of the main army returned.\nAnswer this question: How many of the 30000 men did not survive?

[Answer] 
27500

[Rationale]
The goal of this problem is to calculate how many of the 30000 men did not survive. Before that, we have to check how many men actually survived. In the passage, it is said that "only 2500 remained alive and were captured". Therefore we could conclude that 30000 - 2500 = 27500 men did not survive.



[Example 7]
[Instruction and Question]
The Battle of Stiklestad  in 1030 is one of the most famous battles in the history of Norway. In this battle, King Olaf II of Norway  was killed. During the pontificate of Pope Alexander III, the Roman Catholic Church  declared Olaf a saint in 1164. His younger half-brother, Harald Hardrada, was also present at the battle. Harald was only fifteen when the battle of Stiklestad took place. He became King of Norway in 1047, only to die in a failed invasion of England at the Battle of Stamford Bridge in 1066. The authenticity of the battle as a historical event is subject to question. Contemporary sources say the king was murdered. According to the Anglo Saxon Chronicle of 1030, Olaf was killed by his own people. Adam of Bremen wrote in 1070 that Olaf was killed in an ambush, and so did Florence of Worcester in 1100. Those are the only contemporary sources that mention the death of the king. After the king's canonization it was felt that the saint could not have died in such circumstances.  The story of the Battle of Stiklestad as we know it gradually developed during the two centuries following the death of King Olaf. Saint Olaf must have fallen in a major battle for Christianity.\n\nBased on the above article, answer a question. How many years after Olaf II's death was he declared a saint?

[Answer]
134

[Rationale]
Olaf II of Norway was declared a saint by the Roman Catholic Church during the pontificate of Pope Alexander III in 1164, which is 134 years after his death in the Battle of Stiklestad in 1030.



[Example 8]
    \end{lstlisting}
    




\begin{lstlisting}[caption={Demonstration used for exqa prompt}, label=exqa prompt,,captionpos=b]
[Example 1]
[Instruction and Question]
Extract the answer to the question from the following context.
Question: Architects have to not only be talented designers, but at the other end of the spectrum, have a talent for the more mundane what?
Context: Architecture has to do with planning and designing form, space and ambience to reflect functional, technical, social, environmental and aesthetic considerations. It requires the creative manipulation and coordination of materials and technology, and of light and shadow. Often, conflicting requirements must be resolved. The practice of Architecture also encompasses the pragmatic aspects of realizing buildings and structures, including scheduling, cost estimation and construction administration. Documentation produced by architects, typically drawings, plans and technical specifications, defines the structure and/or behavior of a building or other kind of system that is to be or has been constructed.

[Options]
Not Given

[Answer]
pragmatic aspects of realizing buildings and structures

[Rationale]
The context states architects must have a talent for in addition to design. This includes activities such as scheduling, cost estimation and construction administration which are essential to creating a physical structure. 

Which is why the answer extracted from the context is "pragmatic aspects of realizing buildings and structures".



[Example 2]
[Instruction and Question]
Please answer correctly the following question related to the paragraph below. 
Which tribe was less affected by the disruption of ocean food webs?

Two related native tribes live thousands of miles away from each other. Ana tribe lives near the South Pole in Antarctica, while Momu tribe lives in Northern Alaska near the North Pole.

Hint: Unfortunately, the layer of good ozone is being destroyed by air pollution. The chief culprits are chlorine and bromine gases. They are released in aerosol sprays, coolants, and other products. Loss of ozone has created an ozone hole over Antarctica. Ozone depletion results in higher levels of UV radiation reaching Earth. In humans, this increases skin cancers and eye cataracts. It also disturbs the nitrogen cycle, kills plankton, and disrupts ocean food webs. The total loss of the ozone layer would be devastating to most life. It’s rate of loss has slowed with restrictions on pollutants, but it is still at risk.

[Options]
Not Given

[Answer]
Momu

[Rationale]
The paragraph mentions the negative impacts of the loss of the ozone layer, specifically the disruption of ocean food webs. The Ana tribe lives near the South Pole in Antarctica, where the ozone hole is located, while the Momu tribe lives near the North Pole.

Thus, based on the context of the ozone depletion affecting the ocean food webs, it can be inferred that the Momu tribe was less affected by the disruption of ocean food webs. 

Therefore, the answer is Momu.



[Example 3]
[Instruction and Question]
Given the below context:

Feeling unsettled, Etty left Rome for Venice, intending to remain there for 10 days and then return to England. Evans preferred to remain in Rome, so Etty travelled alone, pausing briefly in Florence and in Ferrara (where he stopped to kiss the armchair of Ludovico Ariosto). The painter Charles Lock Eastlake, then resident in Rome, had provided Etty with a letter of introduction to Harry D'Orville, British Vice consul in Venice; D'Orville was so impressed with Etty that he arranged for him to stay in his own house, rather than in lodgings. Etty had long considered Venice his spiritual home and "The hope and idol of my professional life", and had often wondered why, given its artistic importance, so few English travellers visited the city. He was not disappointed. Throughout the remainder of his life, he looked back on his visit to Venice with great fondness, writing shortly before his death that "Venezia, cara Venezia! thy pictured glories haunt my fancy now!"
Although Etty had only intended to stay for 10 days, he was so taken with Venice that he remained for over seven months. He fell into a routine of copying paintings in Venetian collections by day, and attending the life class of the Venetian Academy of Fine Arts by night, producing around 50 oil paintings in total as well as numerous pencil sketches. He was extremely impressed with the high quality of the Venetian Academy; the instructors in their turn were extremely impressed with the quality of Etty's work, in particular his flesh tones. He acquired the nickname of ""Il Diavolo"" owing to the high speed at which he was able to paint, and watching him at work became something of a spectacle in its own right; luminaries including Gioachino Rossini and Ladislaus Pyrker (then Patriarch of Venice) came to watch him paint. 
So devoted was Etty to his studies in Venice that he exhibited no original work in 1823, writing to his brother that "If one spent all the time in painting originals, one might as well, nay better, be at home". The members of the Venetian...

Guess a valid title for it!

[Options]
Not Given

[Answer]
William Etty

[Rationale]
The context is a description of the experiences and events surrounding the British painter William Etty during his visit to Venice. Throughout the text, the focus is on Etty and his actions, thoughts, and feelings, making it clear that he is the main subject of the paragraph. The text mentions his nickname ""Il Diavolo"" and the admiration he received from notable figures like Gioachino Rossini and Ladislaus Pyrker, further emphasizing his central role.

Given the subject matter and focus on Etty, a fitting title for the paragraph could be "William Etty".



[Example 4]
[Instruction and Question]
I want to test the ability of students to read a passage and answer questions about it. Could you please come up with a good question for the passage?
"Some high-speed black-and-white films, such as Ilford Delta 3200 and Kodak T-MAX P3200, are marketed with film speeds in excess of their true ISO speed as determined using the ISO testing method. For example, the Ilford product is actually an ISO 1000 film, according to its data sheet. The manufacturers do not indicate that the 3200 number is an ISO rating on their packaging. Kodak and Fuji also marketed E6 films designed for pushing (hence the "P" prefix), such as Ektachrome P800/1600 and Fujichrome P1600, both with a base speed of ISO 400."

[Options]
Not Given

[Answer]
What is notable about these films?

[Rationale]
The key point that the passage emphasizes is the difference between the marketed film speeds and the actual ISO speeds of these films.
The passage raises questions about the marketing practices in the photography industry and the reliability of the information provided to consumers.The manufacturers of these films do not indicate the actual ISO ratings on the packaging, which makes it notable. The passage also mentions that these films are marketed for pushing, which is another notable aspect of these films.
Thus, the best answer would be: What is notable about these films?"



[Example 5]
[Instruction and Question]
I am a movie director and I just received the following movie plot. 
Could you help me answer this question? If not, let me know by writing "Not answerable".

Plot title: Bananas
Movie plot: Fielding Mellish (Woody Allen) is the protagonist, but he does not appear until after the opening credits. The cold open, which featured the assassination of the president of the fictional ""banana republic"" of San Marcos and a coup d'état that brings Gen. Emilio Molina Vargas (Carlos Montalban) to power, sets up the situation in which Mellish would become embroiled. The scene was in the form of a championship boxing telecast on Wide World of Sports, with Don Dunphy as the host and Howard Cosell as the commentator.[4] Mellish is a neurotic blue collar man who tries to impress social activist Nancy (Louise Lasser) by trying to get in touch with the revolution in San Marcos. He visits the republic and attempts to show his concern for the native people. However, he is nearly killed by the local caudillo and then saved by the revolutionaries, putting him in their debt. Mellish then learns, clumsily, how to be a revolutionary. When the revolution is successful, the Castro-style leader goes mad, forcing the rebels to place Mellish as their President. When traveling back to the U.S. to obtain financial aid, he reunites with his activist ex-girlfriend and is exposed. In a classic courtroom scene, Mellish tries to defend himself from a series of incriminating witnesses, including a reigning Miss America and a middle-aged African-American woman claiming to be J. Edgar Hoover in disguise. One of the witnesses does provide testimony favorable to Mellish, but the court clerk, when asked to read back this testimony, replies with an entirely different, wholly unfavorable rendition. 
Mellish is eventually sentenced to prison, but his sentence is suspended on the condition that he does not move into the judge's neighborhood. Nancy then agrees to marry him. The film ends with the between-the-covers consummation of their marriage, an event that was over much more quickly than Nancy had anticipated. Like the opening scene, it was accompanied by Cosell providing commentary.

My question: What is Feilding Mellish's occupation?"

[Options]
Not Given

[Answer]
Blue collar worker who becomes president of San Marcos

[Rationale]
In the movie plot, Feilding Mellish is described as a "neurotic blue collar man". The term "blue collar" refers to a type of job that is typically manual labor or involves working with machinery. Blue collar workers often perform physical labor, as opposed to white collar workers who perform administrative, professional, or managerial work. 
So, based on this description, we can conclude that Feilding Mellish's occupation at the start of the movie is a blue-collar worker. Later in the movie, Feilding Mellish becomes involved in a revolution in the fictional banana republic of San Marcos and is eventually placed as its president after the previous leader goes mad.
Therefore, the answer to the question is: "Blue collar worker who becomes President of San Marcos".



[Example 6]
[Instruction and Question]
This article: 
Shackleton's February 1907 announcement that he intended to base his expedition at the old Discovery headquarters was noted by Scott, whose own future Antarctic plans were at that stage unannounced. In a letter to Shackleton, Scott claimed priority rights to McMurdo Sound. "I feel I have a sort of right to my own field of work," he wrote, adding: "anyone who has had to do with exploration will regard this region primarily as mine".He concluded by reminding Shackleton of his duty of loyalty towards his former commander.
Shackleton's initial reply was accommodating: "I would like to fall in with your views as far as possible without creating a position that would be untenable to myself". Edward Wilson, asked by Shackleton to mediate, took an even tougher line than Scott. 
"I think you should retire from McMurdo Sound", he wrote, advising Shackleton not to make any plans to work from anywhere in the entire Ross Sea quarter until Scott decided 
""what limits he puts on his own rights"". 
To this Shackleton replied: "There is no doubt in my mind that his rights end at the base he asked for [...] I consider I have reached my limit and I go no further".The matter was unresolved when Scott returned from sea duty in May 1907. Scott pressed for a line of demarcation at 170° W—everything to the west of that line, including Ross Island, McMurdo Sound, and Victoria Land, would be Scott's preserve. Shackleton, with other concerns pressing on him, felt obliged to concede. 
On 17 May he signed a declaration stating that ""I am leaving the McMurdo base to you"", and that he would seek to land further east, either at the Barrier Inlet visited briefly during the Discovery Expedition, or at King Edward VII Land. He would not touch the coast of Victoria Land at all. It was a capitulation to Scott and Wilson, and meant forfeiting the expedition's aim of reaching the South Magnetic Pole which was located within Victoria Land. 
Polar historian Beau Riffenburgh believes this was "a promise that should never ethically have been demanded and... 

contains an answer for the question: 
"What are the specific names of the two men between whom the dispute soured relations?", what is it ?

[Options]
Not Given

[Answer]
Shackleton

[Rationale]
From the given context, the answer obviously seems to be: Shackleton



[Example 7]
[Instruction and Question]

    \end{lstlisting}
    




\begin{lstlisting}[caption={Demonstration used for mcqa prompt}, label=mcqa prompt,,captionpos=b]
[Example 1]
[Instruction and Question]
Skylar had stopped breathing but Lee held Skylar's arms anyway, in vain. Given the context: Lee want to do what next?
Possible answers: beg the doctors to try again, get her up and walking home, bring Skylar back

[Options]
- beg the doctors to try again
- get her up and walking home
- bring Skylar back

[Answer]
beg the doctors to try again

[Rationale]
The context of the situation is that Skylar has stopped breathing and Lee is holding Skylar's arms. This suggests that Skylar has received medical attention and is no longer alive.
Based on this, Lee would want to beg the doctors to try again, as this would be the best way to potentially save Skylar's life. 
The options of getting Skylar up and walking home or bringing Skylar back do not align with the context, as they would not address the immediate medical emergency at hand.

So the answer is to beg the doctors to try again.



[Example 2]
[Instruction and Question]
Do you think the right answer to the question "what can run on alcoholic fermentation of the glucose in corn or other plants?" is "bicycles"?

[Options]
- No
- Yes

[Answer]
No

[Rationale]
Alcoholic fermentation is a process that occurs in living organisms, most commonly yeast, to convert glucose into energy. This process does not produce enough energy to power something like a bicycle.

So the answer is No.



[Example 3]
[Instruction and Question]
Eventually , the cyst can get so large that part of the wall weakens and fluid begins to leak out ; that is , the cyst ruptures . It is similar to popping a water balloon . When this happens , the fluid may irritate the pelvic lining and cause pain , often severe . Most of the time the pain begins to gradually subside over the course of the next few days .
What type of doctor might the person have had to see ?

- A paediatrician  - A heart surgeon - a gynecologist - None of the above choices

[Options]
- A paediatrician 
- A heart surgeon 
- a gynecologist 
- None of the above choices

[Answer]
a gynecologist

[Rationale]
The context mentions a cyst in the pelvic region and the symptoms of the cyst rupturing, including pain and fluid leakage. This suggests that the issue is related to the reproductive or urinary system. Given that the issue is related to the reproductive or urinary system, it is likely that the person would have had to see a specialist in these areas. 
Out of the given options, "a gynecologist" is the most likely specialist that the person would have had to see, as gynecologists specialize in female reproductive health, including issues related to the uterus, ovaries, and other reproductive organs. 

In conclusion, based on the given context, the most likely answer is a gynecologist.



[Example 4]
[Instruction and Question]
Process:
- Get some seeds
- Pick a spot to plant them
- Dig a hole in the dirt
- Put the seed in the hole
- Pour some water on the seed and hole
- Cover up the hole
- Press down on it
- Spray some plant food on it.

Question:
suppose if no water is available happens, how will it affect vegetables?

- A: more
- B: less

[Options]
- A
- B

[Answer]
B

[Rationale]
The process of planting seeds requires water, which is essential for the growth of plants. Water helps transport nutrients, regulate temperature, and absorb minerals from the soil. If no water is available, the seed will not germinate and the plant will not grow properly, leading to fewer vegetables. Hence, the absence of water will result in "less" vegetables, making the answer "B".



[Example 5]
[Instruction and Question]
Read the following conversation and answer the question.
M: I'm worried about my sister, Lucy.
W: Why is that?
M: She thinks she's too fat.
W: Is she?
M: No, but she keeps skipping meals. Then, she only eats chips and drinks cola.
W: I used to do that. It's called binge. It was no fun!
M: Why did you stop doing it?
W: Well, my doctor told me to eat when I'm hungry. She said, "Eat till you're full or you'll eat too much later." She said a lot of girls ruin their health this way.
M: Did she say what to eat?
W: She said, "Eat fruit, vegetables, meats and grains. Have regular meals and snacks. Get exercise, too."

Question: What's the relationship between the speakers? 

- Doctor and patient.
- Friends.
- Brother and sister.

[Options]
- Doctor and patient
- Friends
- Brother and sister

[Answer]
Friends

[Rationale]
The conversation between two people, a man(M) and a woman(W), is discussing the man's sister who has a negative view of her own body and is skipping meals. The woman shares that she used to engage in similar behavior, and explains that her doctor advised her to eat a balanced diet and exercise.

Among the three options,
- First, since the woman is sharing her experience of what the doctor told her, it is unlikely she would be a medical professional. So the answer couldn't be doctor and patient.
- Second, given that the man is talking about his sister, the woman is unlikely to be the man's sister.
- Lastly, because the conversation is informal and the speakers are discussing a personal matter it is likely that the relationship between the speakers is that of friends.

Therefore, the answer is Friends.



[Example 6]
[Instruction and Question]
Question: The best place for the troublemaker to sneak a smoke while he was in where, was in the bathroom?

Choices: 
- playground
- school
- at hotel
- theater
- convenience store

The rationale to choose ""school"" as the answer is that:

[Options]
Not Given

[Answer]
smokers traditionally use school bathrooms

[Rationale]
Here's the rationale: "smokers traditionally use school bathrooms".



[Example 7]
[Instruction and Question]
So I was looking at quotes for our tshirts for ice hockey , and I came across one that sums up how I ' m feeling right now . When you talk about an injury and the kind of depression you go through , it 's not just because you 're out of shape and you can't go out and play . You 're missing a part of you .

Why is the narrator so down in the dumps ?
Pick the best answer from the following options:
A. They are unable to play ice curling .
B. They are hurt .
C. They are unable to play ice skating .
D. They are unable to play in the yard .

[Options]
- A
- B
- C
- D

[Answer]
B

[Rationale]
The quote in the question mentions the narrator's feelings about an injury and the depression that follows. The narrator says that it's not just about being out of shape and unable to play, but rather they feel like they are missing a part of themselves. This implies that the reason for their depression is linked to their physical inability to participate in a sport or activity they enjoy. Based on the context, it can be inferred that the narrator is upset because they are hurt and unable to play ice hockey.
So, the answer is B.



[Example 8]
[Instruction and Question]
Q: What does the ubiquitous generator use to generate currents?

Choices:

- inertia - water - gravity - magnetism

[Options]
- inertia
- water
- gravity
- magnetism

[Answer]
magnetism

[Rationale]
A generator is a device that converts mechanical energy into electrical energy. This is done by using a magnetic field to induce a current in the generator’s windings. This current is then sent to an external circuit, where it can be used to power electrical devices. The generator uses magnetism to generate currents, so the answer is magnetism.



[Example 9]
[Instruction and Question]

    \end{lstlisting}



\begin{lstlisting}[caption={Demonstration used for sni commonsense}, label=sni commonsense,,captionpos=b]
[Example 1]
[Instruction and Question]
In this task, you will be shown a short story with a beginning, two potential middles, and an ending. Your job is to choose the middle statement that makes the story incoherent / implausible by indicating 1 or 2 in the output. If both sentences are plausible, pick the one that makes less sense.\n\nBeginning: May needed to lose a lot of weight. Middle 1: she decided to hire a fitness trainer and dietician. Middle 2: she decided to hire a fitness trainer and dietician, but it didn't help. Ending: With their help, May lost forty pounds!

[Answer]
2

[Rationale]
Middle 1 indicates that May took action to achieve her goal of losing weight by hiring a fitness trainer and dietician, and with their help, she succeeded in losing forty pounds. Middle 2 also indicates that May hired a fitness trainer and dietician, but despite their help, she didn't achieve her goal of losing weight. So Middle 2 makes less sense.



[Example 2]
[Instruction and Question]
In this task, you are given two phrases: Head and Tail, separated with <sep>. The Head and the Tail events are short phrases possibly involving participants. The names of specific people have been replaced by generic words (e.g., PersonX, PersonY, PersonZ). PersonX is always the subject of the event. You have to determine whether The Tail is the reason for the Head or not. The reason provides a post-fact explanation of the cause of an event. For example, why one has to walk could be explained by a car has broken down. Classify your answers into \"Yes\" and \"No\". The phrase may also contain \"___\", a placeholder that can be an object, a person, and/or an action.\n\nHead: PersonX accepts the invitation<sep>Tail: makes new friends

[Answer]
No

[Rationale]
While making new friends could be a potential outcome of accepting the invitation, it does not explain why PersonX accepted the invitation in the first place. Therefore the Tail does not directly provide a reason for the Head.



[Example 3]
[Instruction and Question]
You will be given a passage consisting of set of facts and a question as input. The task is to answer a question of form 'Where is <person_name>?' using one of the given facts to determine the latest location of the person. Answer should be a word/phrase describing the location from the supporting fact. Avoid answers that are incomplete or incorrect.\n\nPassage: Daniel went to the bedroom. Mary went to the kitchen. Mary travelled to the bedroom. John journeyed to the bathroom. Sandra moved to the garden. Sandra went to the office. Daniel went back to the hallway. Sandra moved to the kitchen. Question: Where is Sandra?

[Answer]
kitchen

[Rationale]
We can be infer from the last fact in the passage which states "Sandra moved to the kitchen."



[Example 4]
[Instruction and Question]
Two analogies that relate places/locations to the associated travel mode is given in the form \"A : B. C : ?\". \"A : B\" relates place A to travel mode B. Your task is to replace the question mark (?) with the appropriate travel mode for the given place C, following the \"A : B\" relation.\n\nhotel : taxi. france : ?

[Answer]
airplane

[Rationale]
To go to the hotel, one could take a taxi. Similarily, one could go to France by taking an airplane.



[Example 5]
[Instruction and Question]
Given an object and a part, decide whether the object has that part. For example is you are asked 'gun has barrel', you need to decide if a gun has a barrel as one of its components or parts. All sentences strictly follow the template 'object has part?.' The answer should be 1 or 0, 1 means the object has the given part, while 0 means it doesn't have the part.\n\nhusband has wife?

[Answer]
1

[Rationale]
A husband is a term used to describe a married man, which implies that he has a wife as his partner in marriage. Therefore, a husband has a wife.



[Example 6]
[Instruction and Question]
You are provided with an \"Event\" and it's \"Intent\" related to PersonX. Determine the sentiment value of the given input as either \"Positive\", \"Negative\", and \"Unknown\". \n\nEvent:PersonX talks to PersonY about anything. Intent: 1) protect themselves 2) stand up for themselves

[Answer]
Unknown

[Rationale]
The intent of "protect themselves" and "stand up for themselves" could potentially suggest a negative or positive sentiment depending on the context and the nature of the conversation between PersonX and PersonY. Additionally, the event of simply "talking about anything" is too vague to infer any sentiment.



[Example 7]
[Instruction and Question]
In this task, you will be given a short story. One sentence from the story is chosen. Consider the likely emotions and basic human drives of the participants in that sentence. Does any of these states of mind/feelings motivate the participant to do what happens in that sentence? You should write your answer in the form \" A >Motivates> B\". Try to use phrases and sentences from the story to compose your answer when possible. For the motivation sentence, you must choose a verb from :feel(s), want(s) or like(s). There will always be some motivation in the given story.\n\nstory: Ed took his dog Molly to the dog park. He was feeling sad because he and his girl had recently broken up. As he sat watching his dog play, a pretty woman sat next to him. They started chatting and really hit it off. They made a date for the following weekend for dinner.\n selected sentence: As he sat watching his dog play, a pretty woman sat next to him.

[Answer]
Ed like(s) watching his dog play >Motivates> Ed watches his dog play

[Rationale]
Ed's enjoyment of watching his dog play motivates him to sit and watch his dog at the dog park.



[Example 8]
[Instruction and Question]
In this task, we ask you to write an answer to a question that involves event \"frequency\", which refers to how often an event is likely to be repeated. For example, \"taking showers\" typically occurs ~5 times a week, \"going to Saturday market\" usually happens every few weeks/months, etc. Note that a lot of the questions could have more than one correct answer. We only need a single most-likely answer. Please try to keep your \"answer\" as simple as possible. Concise and simple \"answer\" is preferred over those complex and verbose ones.\n\nSentence: Tommy and Suzy (brother and sister) went to the playground one afternoon with their mom and dad, Jan and Dean. \nQuestion: How often did they go to the playground?

[Answer]
three times a week

[Rationale]
The family would go to the playground on a regular basis, specifically three times a week.



[Example 9]
[Instruction and Question]
Given a sentence with a missing word, pick the answer option that best fills out the missing word in the sentence. Indicate each answer with its index ('a', 'b', 'c', 'd').\n\nUnconsciousness can take place within one to ____ minutes.\\Question: Choose the right answer from options given a) no b) two c) one d) zero

[Answer]
b

[Rationale]
The sentence is discussing how quickly unconsciousness can occur. Before choosing among the options, it is good to remind there is a phrase "one to two minutes" meaning that something happens very quickly.



[Example 10]
[Instruction and Question]
In this task, you're given four sentences of a story written in natural language in which one part is missing. Your job is to predict the position and missing part of the story and return in the following format: position, missing part. The missing part is a sentence that completes the story, and the position is the number of the missing sentence in the new story.\n\nSentence1: Fred always stopped at the same coffee shop before work. Sentence2: He got to know the attractive barista. Sentence3: He developed a crush on her. Sentence4: She said yes.

[Answer]
4, Finally Fred worked up the courage to ask her on a date

[Rationale]
The story does not explain why the barista said yes to Fred. It is most likely that Fred asked her to go on a date with him. So the completed story would look like:\nSentence1: Fred always stopped at the same coffee shop before work.\nSentence2: He got to know the attractive barista.\nSentence3: He developed a crush on her.\nMissing Sentence: Finally Fred worked up the courage to ask her on a date.\nSentence4: She said yes.\n\nSince the missing sentence is inserted in front of Sentence4, the first answer is "4" and the second answer is "Finally Fred worked up the courage to ask her on a date".



[Example 11]
[Instruction and Question]
In the following task, you are given a yes/no question, its answer, and additional information which includes a description of a key term in the question and several facts. Your task is to \"decompose the question\"i.e. write the steps required to construct the given answer where each step is a question that can be answered using Wikipedia articles. To decompose the question, think of a decomposition strategy you could apply to answer the question through a search against Wikipedia. The decomposition strategy might not be obvious for some questions. For this reason, you are provided with additional information. You are not required to write decompositions that exactly match the provided facts, they are only meant as guidance. If the facts are too noisy, think about how you would have used Wikipedia to answer this question to create the decomposition. \"not_definitive_answer\" is true if there is no clear answer i.e. the question is ambiguous or a subjective opinion.\"incorrect_answer\" is true if the answer to the question is not correct. If the answer is incorrect, write a decomposition that evaluates the correct answer. Each question should be easy to answer based on simple facts from Wikipedia pages, or a simple check using the answers to preceding questions. To ensure this, we also ask you to enter a potential Wikipedia article that would contain the answer to each question. If you can not think of a Wikipedia page that can answer your question, please think of a different strategy. You do not need to enter a Wikipedia page for steps that don't add new information and only contain simple checks. The answer to the final question is the same as the answer to the original question. Every question's answer has an impact on the final answer and there should be no redundant steps in the decomposition.\n\nQuestion: Can oysters be preserved without refrigeration?, Answer:Yes, Description of Oyster: salt-water bivalve mollusc, Fact1: In some types of Korean kimchi, oysters are placed between the leaves of nappa cabbage. , Fact2: Many grocery stores carry canned oysters in the shelf stable section. 

[Answer]
Step1: How are oysters preserved in various types of Korean kimchi?, Step2: What are the common methods of preserving oysters in grocery stores?, Step3: Do any of #1 or #2 not require refrigeration?, Wikipedia page for step 1: Kimchi, Wikipedia page for step 2: Oyster, Wikipedia page for step 3: , not_definitive_answer: False, incorrect_answer: False

[Rationale]
The question is asking if oysters can be preserved without refrigeration. To answer this, we need to know how oysters are preserved in general. The description of oysters tells us that they are a salt-water bivalve mollusc.\nTo decompose the question, we can start by asking how oysters are preserved in various types of Korean kimchi (Step1), which would lead us to the Wikipedia page for Kimchi. Next, we can ask what the common methods of preserving oysters in grocery stores are (Step2), which would lead us to the Wikipedia page for Oyster.\nBased on the information we gather from Step1 and Step2, we can ask if any of the methods described do not require refrigeration (Step3). If we find that some methods do not require refrigeration, we can conclude that oysters can be preserved without refrigeration.



[Example 12]
[Instruction and Question]
Given a premise, an initial context, an original ending, and a counterfactual context, the task is to generate a new story ending aligned with the counterfactual context and as close to the original ending as possible. Each instance consists of a five-sentence story. The premise is the first sentence of a story, and the second sentence, which is the initial context, provides more information about the story's context and the story's general plot. The original ending is the last three sentences of the story. Also, a counterfactual context is a slight modification to the initial context. You should write a new story ending that edits the original story ending as little as possible to regain coherence with the counterfactual context. To sum up, you should write the last three sentences of a story based on the premise(first sentence) and the counterfactual context(second sentence) of the story.\n\nPremise: I used to pitch for my high school baseball team. \n Initial Context: I was a very good pitcher. \n Original Ending: I spent my weekends throwing fastballs with my dad. All that practice really helped me out. I feel like I had developed a perfect fastball. \n Counterfactual Context: I was so bad they moved me to right field.

[Answer]
I spent the weekends catching fly balls with my dad. All that practice really helped me out. I feel like I had developed into a great right fielder.

[Rationale]
The original ending implies that the protagonist is a successful pitcher, but the counterfactual context changes that by stating that they were moved to right field due to their lack of skill. To regain coherence with the counterfactual context, the new ending needs to reflect this change. Instead of practicing fastballs, the protagonist now spends their weekends catching fly balls with their dad, which is a suitable activity for a right fielder. The rest of the new ending is similar to the original one, with the protagonist feeling like they have improved their skills.



[Example 13]
[Instruction and Question]

    \end{lstlisting}
    




\begin{lstlisting}[caption={Demonstration used for sni formallogic}, label=sni formallogic,,captionpos=b]
[Example 1]
[Instruction and Question]
This task evaluates for the ability to follow basic natural language instructions nested and performing a sequence of operations, including basic logic and conditionals.\n\nrepeat a woodchuck chucks lots of wood two times, but replace lots with five pounds the first time and two tons the second time

[Answer]
a woodchuck chucks five pounds of wood a woodchuck chucks two tons of wood

[Rationale]
Step1: repeat a woodchuck chucks lots of wood two times, Intermediate Answer 1: a woodchuck chucks lots of wood a woodchuck chucks lots of wood\nStep 2: replace lots with five pounds the first time and two tons the second time, Final Answer: a woodchunk chunks five pounds a woodchunk chunks two tons of wood"



[Example 2]
[Instruction and Question]
In this task, you are given a year. You need to check if it is a leap year or not. A year may be a leap year if it is evenly divisible by 4. Years that are divisible by 100 (century years such as 1900 or 2000) cannot be leap years unless they are also divisible by 400. Return 1 if it is a leap year, else return 0.\n\n1244

[Answer]
1

[Rationale]
The year 1244 is divisible by 4 with no remainder, yet not divisible by 100. Therefore, 1244 is a leap year and the answer should be 1.



[Example 3]
[Instruction and Question]
In this task, you are given a date in \"dd/mm/yyyy\" format. You need to check if the date is valid or not. Return 1 if it is valid, else return 0. A date is valid if the components day(\"dd\"), month(\"mm\") and year(\"yyyy\") are all valid individually. A day(dd) is valid if it: a) lies between 1 and 31 for the months of January, March, May, July, August, October, December, b) lies between 1 and 28 for February on non-leap years and between 1-29 on leap years, c) lies between 1-30 for all other months. A month(mm) is valid if it lies in the range from 1 to 12 as there are 12 months in a year. A year is always valid if it is expressed in the form of \"yyyy\".\n\n31/18/1191

[Answer]
0

[Rationale]
The given date "31/18/1191" is not a valid date. The month component is "18" which is not in the range of 1-12. Therefore, the date is invalid and the output should be 0.



[Example 4]
[Instruction and Question]
In this task, you are given a date in a particular format and you need to convert to another format. If given format is \"dd/mm/yyyy\" then convert to \"mm/dd/yyyy\". If given format is \"mm/dd/yyyy\" then convert to \"dd/mm/yyyy\".\n\n04/17/1955, input_format=mm/dd/yyyy

[Answer]
17/04/1955

[Rationale]
Since there is no 17th month, the given format is "mm/dd/yyyy". We have to convert it to "dd/mm/yyyy", so the answer is "17/04/1995".



[Example 5]
[Instruction and Question]
You are given a time in 24-Hours format, and you need to convert it to time in the 12-Hours format. For a 24-Hours format time larger than 12:00, subtract 12 hours from the given time, then add 'PM'.  For example, if you have 14:30 hours, subtract 12 hours, and the result is 2:30 PM. If the 24-Hours format time is less than or equal to 12:00, add 'AM'. For example, say you have 10:15 hours, add the 'AM' to the end, here we get 10:15 AM. Note that 00:00 Hrs in 24-Hours format is 12:00 AM in 12-Hours format and 12:00 Hrs in 24-Hours format would be 12:00 PM in 12-Hours format.\n\n21:50 Hrs

[Answer]
9:50 PM

[Rationale]
The given time is 21:50 Hrs in the 24-Hours format.\nSince it is greater than 12:00, we need to subtract 12 hours to convert it to the 12-Hours format.\n\n21:50 - 12:00 = 9:50\nSo, the time in the 12-Hours format is 9:50 PM.



[Example 6]
[Instruction and Question]
Given a statement about date and time, state whether the statement is true or false. The number of date/time operands in the statement ranges between 2 and 3. Let's say the values are denoted by t1, t2 and t3. The statements follow one of the following ten templates: 't1 occurs before t2, t1 doesn't occur before t2, t1 occurs after t2, t1 doesn't occur after t2, t1 occurs between t2 and t3, t1 doesn't occur between t2 and t3, t1 occured before t2 but after t3, t1 occured after t2 but before t3, t1 didn't occur before t2 but after t3, t1 didn't occur after t2 but before t3'. The output should be either 'True' or 'False'.\n\n12:38:06 occurs between 0:56:08 and 07:36:45 PM

[Answer]
True

[Rationale]
The given range of time is 0:56:08 to 7:36:45 PM which is in 12-hour format. The given time 12:38:06 is also in 12-hour format. Therefore, we need to convert the given range to 24-hour format, which would be 00:56:08 to 19:36:45. Since 12:38:06 falls between the given range, the statement is true.



[Example 7]
[Instruction and Question]

    \end{lstlisting}
    





\begin{lstlisting}[caption={Demonstration used for sni logic}, label=sni logic,,captionpos=b]
[Example 1]
[Instruction and Question]
This task evaluates for the ability to follow basic natural language instructions nested and performing a sequence of operations, including basic logic and conditionals.\n\nrepeat a woodchuck chucks lots of wood two times, but replace lots with five pounds the first time and two tons the second time

[Answer]
a woodchuck chucks five pounds of wood a woodchuck chucks two tons of wood

[Rationale]
Step1: repeat a woodchuck chucks lots of wood two times, Intermediate Answer 1: a woodchuck chucks lots of wood a woodchuck chucks lots of wood\nStep 2: replace lots with five pounds the first time and two tons the second time, Final Answer: a woodchunk chunks five pounds a woodchunk chunks two tons of wood"



[Example 2]
[Instruction and Question]
In this task, you are given a year. You need to check if it is a leap year or not. A year may be a leap year if it is evenly divisible by 4. Years that are divisible by 100 (century years such as 1900 or 2000) cannot be leap years unless they are also divisible by 400. Return 1 if it is a leap year, else return 0.\n\n1244

[Answer]
1

[Rationale]
The year 1244 is divisible by 4 with no remainder, yet not divisible by 100. Therefore, 1244 is a leap year and the answer should be 1.



[Example 3]
[Instruction and Question]
In this task, you are given a date in \"dd/mm/yyyy\" format. You need to check if the date is valid or not. Return 1 if it is valid, else return 0. A date is valid if the components day(\"dd\"), month(\"mm\") and year(\"yyyy\") are all valid individually. A day(dd) is valid if it: a) lies between 1 and 31 for the months of January, March, May, July, August, October, December, b) lies between 1 and 28 for February on non-leap years and between 1-29 on leap years, c) lies between 1-30 for all other months. A month(mm) is valid if it lies in the range from 1 to 12 as there are 12 months in a year. A year is always valid if it is expressed in the form of \"yyyy\".\n\n31/18/1191

[Answer]
0

[Rationale]
The given date "31/18/1191" is not a valid date. The month component is "18" which is not in the range of 1-12. Therefore, the date is invalid and the output should be 0.



[Example 4]
[Instruction and Question]
In this task, you are given a date in a particular format and you need to convert to another format. If given format is \"dd/mm/yyyy\" then convert to \"mm/dd/yyyy\". If given format is \"mm/dd/yyyy\" then convert to \"dd/mm/yyyy\".\n\n04/17/1955, input_format=mm/dd/yyyy

[Answer]
17/04/1955

[Rationale]
Since there is no 17th month, the given format is "mm/dd/yyyy". We have to convert it to "dd/mm/yyyy", so the answer is "17/04/1995".



[Example 5]
[Instruction and Question]
You are given a time in 24-Hours format, and you need to convert it to time in the 12-Hours format. For a 24-Hours format time larger than 12:00, subtract 12 hours from the given time, then add 'PM'.  For example, if you have 14:30 hours, subtract 12 hours, and the result is 2:30 PM. If the 24-Hours format time is less than or equal to 12:00, add 'AM'. For example, say you have 10:15 hours, add the 'AM' to the end, here we get 10:15 AM. Note that 00:00 Hrs in 24-Hours format is 12:00 AM in 12-Hours format and 12:00 Hrs in 24-Hours format would be 12:00 PM in 12-Hours format.\n\n21:50 Hrs

[Answer]
9:50 PM

[Rationale]
The given time is 21:50 Hrs in the 24-Hours format.\nSince it is greater than 12:00, we need to subtract 12 hours to convert it to the 12-Hours format.\n\n21:50 - 12:00 = 9:50\nSo, the time in the 12-Hours format is 9:50 PM.



[Example 6]
[Instruction and Question]
Given a statement about date and time, state whether the statement is true or false. The number of date/time operands in the statement ranges between 2 and 3. Let's say the values are denoted by t1, t2 and t3. The statements follow one of the following ten templates: 't1 occurs before t2, t1 doesn't occur before t2, t1 occurs after t2, t1 doesn't occur after t2, t1 occurs between t2 and t3, t1 doesn't occur between t2 and t3, t1 occured before t2 but after t3, t1 occured after t2 but before t3, t1 didn't occur before t2 but after t3, t1 didn't occur after t2 but before t3'. The output should be either 'True' or 'False'.\n\n12:38:06 occurs between 0:56:08 and 07:36:45 PM

[Answer]
True

[Rationale]
The given range of time is 0:56:08 to 7:36:45 PM which is in 12-hour format. The given time 12:38:06 is also in 12-hour format. Therefore, we need to convert the given range to 24-hour format, which would be 00:56:08 to 19:36:45. Since 12:38:06 falls between the given range, the statement is true.



[Example 7]
    \end{lstlisting}
    




\begin{lstlisting}[caption={Demonstration used for sni program execution}, label=sni program execution,,captionpos=b]
[Example 1]
[Instruction and Question]
In mathematics, the absolute value of a number is the non-negative value of that number, without regarding its sign. For example, the absolute value of -2 is 2, and the absolute value of 5 is 5. In this task you will be given a list of numbers and you need to return the element with highest absolute value. If a negative and positive element have the same absolute value you should return the positive element. The absolute value for negative numbers can be found by multiplying them by -1. After finding the element with the maximum absolute value you should return the value of that element before you applied the absolute value.\n\n[-36.642 -92.499 -81.859  71.021 -90.55  -98.633 -29.564  58.697  62.083\n -56.353]

[Answer]
-98.633

[Rationale]
The element with the highest absolute value in the given list is -98.633, which has an absolute value of 98.633. The absolute value of all other numbers in the list is less than or equal to 98.633. However, we need to return the value of the element before applying the absolute value, so the answer is -98.633.



[Example 2]
[Instruction and Question]
In this task you are given a list of numbers and you need to find the average of each two consecutive values. The average of two numbers a and b is calculated as: (a + b) /2. The output should be a list of the averages of each two consecutive values. A list is presented with two brackets and comma-separated values, like: [1,2,3].\n\n[-24, -89, -70, 48, -50, 86, 66, 92, 70, 85, -72]

[Answer]
[-56.5, -79.5, -11.0, -1.0, 18.0, 76.0, 79.0, 81.0, 77.5, 6.5]

[Rationale]
The average of -24 and -89 is (-24 + (-89)) / 2 = -56.5\nThe average of -89 and -70 is (-89 + (-70)) / 2 = -79.5\nThe average of -70 and 48 is (-70 + 48) / 2 = -11\nThe average of 48 and -50 is (48 + (-50)) / 2 = -1\nThe average of -50 and 86 is (-50 + 86) / 2 = 18\nThe average of 86 and 66 is (86 + 66) / 2 = 76\nThe average of 66 and 92 is (66 + 92) / 2 = 79\nThe average of 92 and 70 is (92 + 70) / 2 = 81\nThe average of 70 and 85 is (70 + 85) / 2 = 77.5\nThe average of 85 and -72 is (85 + (-72)) / 2 = -17.5



[Example 3]
[Instruction and Question]
In this task, you are given inputs i and A, where i is an integer and A is a list. You need to output the first i elements of A, starting with the 1st element in A. i will always have a value less than the length of A\n\n4, ['1189', 'f', '7341', '7391', 'p', '8555', '5925', '73', 'O', 'q', '2401', 'S', '3803', 'r', 'U']

[Answer]
1189, f, 7341, 7391

[Rationale]
Given the list ['1189', 'f', '7341', '7391', 'p', '8555', '5925', '73', 'O', 'q', '2401', 'S', '3803', 'r', 'U'] as input, we need to find the first 4 elements:\n1. The 1st element is 1189.\n2. The 2nd element is 'f'.\n3. The 3rd element is '7341'.\n4. The 4th element is '7391'.\n\nThe final output would be ['1189', 'f', '7341', '7391'].



[Example 4]
[Instruction and Question]
You are given a list of integers and an integer target, return a list of a pair of numbers in any order such that they add up to target. If there is no such pair of numbers, then return an empty list\n\n[42, 22, 36, 49, 31, 33, 27, 24, 37, 26, 34, 46, 16, 29, 40, 41, 2, 39, 11, 17], target=88

[Answer]
[42, 46]

[Rationale]
Among list, 42 and 46 adds up to 88.



[Example 5]
[Instruction and Question]
In this task, you are given a string of characters. You need to remove duplicate characters from the string if any, and return the resulting string.\n\nPgnuLHMWcjJVzGyJelrh

[Answer]
PgnuLHMWcjJVzGyelrh

[Rationale]
Since there are no duplicate characters in the given string, the resulting string will be the same as the input string. Therefore, the output will be "PgnuLHMWcjVzGyJelrh".



[Example 6]
[Instruction and Question]
In this task, you need to count the number of vowels (letters 'a', 'e', 'i', 'o', 'u') / consonants (all letters other than vowels) in the given sentence.\n\nSentence: 'a metro bus stopped at an intersection with tall building in the background'. Count the number of vowels in the given sentence.

[Answer]
23

[Rationale]
Given the sentence 'a metro bus stopped at an intersection with tall building in the background', let's try one word-by-word.\n1. 'a' : 1 -> (total) 1\n2. 'metro' : 2 -> (total) 3\n3. 'bus' : 1 -> (total) 4\n4. 'stopped' : 2 -> (total) 6\n5. 'at' : 1 -> (total) 7\n6. 'an' : 1 -> (total) 8\n7. 'intersection' : 5 -> (total) 13\n8. 'with' : 1 -> (total) 14\n9. 'tall' : 1 -> (total) 15\n10. 'building' : 3 -> (total) 18\n11. 'in' : 1 -> (total) 19\n12. 'the' : 1 -> (total) 20\n13. 'background' : 3 -> (total) 23\n\nThe final total number is 23.



[Example 7]
[Instruction and Question]
In this task you will be given a string of characters. You should remove all vowels from the given string. Vowels are: i,e,a,u,o. The character 'y' or 'Y' does not count as a vowel.\n\nEDIOEOXxdSzUiKOPZ

[Answer]
DXxdSzKPZ

[Rationale]
The given string contains the vowels 'E', 'I', 'O', 'E', 'O', 'U', and 'I', which should be removed according to the instructions. The resulting string after removing the vowels is 'DXxdSzKPZ'.



[Example 8]
[Instruction and Question]
In this task, you are given inputs i,j, and A, where i and j are integers and A is a list. You need to return the sum of all the numerical elements in the list A between the positions i and j (including positions i and j). Return 0 if no numerical element is present in the list between the given ranges. i and j will be non-negative, and will always have a value less than the length of A. i will always be less than j.\n\n12, 15, ['7933', '6321', '6815', 'r', 'y', '2015', 'I', 'e', '7733', '1969', '1973', '3141', '6607', 'n', '5417', '5927', '8123', 'j', 'G', 'f', 'F', '3305', 'M', 'E', 'v', 'c']

[Answer]
15165

[Rationale]
Given the list ['7933', '6321', '6815', 'r', 'y', '2015', 'I', 'e', '7733', '1969', '1973', '3141', '6607', 'n', '5417', '5927', '8123', 'j', 'G', 'f', 'F', '3305', 'M', 'E', 'v', 'c'], we need to find the elements within the range 12 ~ 15.\n1. The 12nd element is 3141.\n2. The 13th element is 6607.\n3. The 14th element is n. This is converted into 0 during the final calculation.\n4. The 15th element is 5417.\n\nCalculating the final results, 3141 + 6607 + 0 + 5417 = 9748 + 5417 = 15165.



[Example 9]
[Instruction and Question]

    \end{lstlisting}
    




\begin{lstlisting}[caption={Demonstration used for sni toxic}, label=sni toxic,,captionpos=b]
[Example 1]
[Instruction and Question]
Categorize the comment on the basis of obscenity. If the comment is obscene output Yes, otherwise output No.\n\n\"Dams thwart salmon migration, degrade water quality, alter water flows and contribute to fish diseases and algae bloom problems.\"\n\nMeanwhile. the sudden rush of \"experts\" on hydrology, power production, etc. pile up like train wreck online and clearly don't give \"dam\" about salmon migration or water quality.

[Answer]
No

[Rationale]
The given comment is a criticism of the construction of dams and how it affects the environment and wildlife. The second sentence expresses frustration towards people who claim to be experts but do not seem to care about the issues mentioned in the first sentence. However, this frustration is not expressed through any obscenities or offensive language.



[Example 2]
[Instruction and Question]
Given a prompt and two completions from a language model, select the completion that is the most fluent with respect to continuing the prompt. Fluency is defined by coherency and ease of understanding, not necessarily grammatical correctness. Generate either Completion A or Completion B.\n\nPrompt: Thanks for small flat rate priority boxes. I'll remember to use \nCompletion A: this. \nCompletion B: them. (Thanks to David for this.)So, how do we get this thing up to speed on time?Well, i

[Answer]
Completion B

[Rationale]
Completion A is not a complete sentence and does not add any meaningful continuation to the prompt. Completion B is the more fluent and coherent choice because it maintains the grammatical structure of the prompt and refers to "small flat rate priority boxes" using the appropriate pronoun "them"



[Example 3]
[Instruction and Question]
In this task you should classify the given text as 'violence' if there is inhumanity associated with it. Otherwise, classify it as 'not violence'. Mentioning killing people, or hurting them in any way is considered violent.\n\nIma fucking stab u up bruv. Fucking white prick bruv.. fucking racist.  Gammon

[Answer]
violence

[Rationale]
The derogatory language and racial slurs contribute to the inhumanity associated with the text.The use of violent language and threatening to physically harm someone (stabbing) makes this statement violent.



[Example 4]
[Instruction and Question]
Classify the given tweet into the three categories: (1) 'Hate Speech', (2) 'Offensive' and (3) 'Neither'. 'Hate Speech' is kind of a threating statement or sometimes include call for violence while 'offensive' statement just offensds someone. 'Neither' is when it doesn't fall into Hate Speech or Offensive category.\n\n&#128514;&#128514;&#128514;&#128514;&#128514;&#128514; RT @KarltonHinez: @ItsNotHarold smh bitch sellin her ass

[Answer]
Offensive

[Rationale]
The given tweet includes the word "bitch" which is a derogatory term for a woman. Additionally, the tweet refers to someone selling their body, which could be interpreted as objectifying and degrading. While the tweet does not contain a direct call for violence or threaten anyone, it is still offensive and inappropriate.



[Example 5]
[Instruction and Question]
In this task, you are given a hateful post in English from online platforms. You are expected to classify the post into two classes: aggresive or non-aggresive. An aggresive post is one that expresses feelings and opinions in an abusive way and tries to dominate others. Note that the URLs in the text have been replaced with [Link].\n\nPost: 70% of Credible Fear Migrants Caught By Border Patrol Once Released into the U.S they never file Asylum ICE has to find them amd deport     #RedNationRising #Trump #MAGA  Illegal's #StopAsylumScam  Illegal #FamiliesBelongTogether in their Country Not the U.S Stop #immigration [Link]

[Answer]
Non-aggresive

[Rationale]
The given post expresses some political opinions and criticism about immigration policies, but it does not contain any abusive language or attempt to dominate others. Therefore, it can be classified as non-aggressive.



[Example 6]
[Instruction and Question]
The input is a tweet which can be classified as Hate Speech, Offensive or Normal. Given such a tweet, output the class the tweet belongs to. Hate Speech will contain threatening language towards the community targeted. Offensive language will contain abusive or discriminatory language towards the community targeted.\n\nthere are approx <number> <number> jews there are over <number> <number> muslims and those animals are breeding like rabbits i prioritize targets the jew population is stable and low the sandnigger population is growing like plague which one do you think we should go for first

[Answer]
Hate Speech

[Rationale]
The use of terms like "animals" and "breeding like rabbits" to describe the Muslim community is dehumanizing and shows a clear intent to incite hatred and violence. The statement also suggests that the person has prioritized targets based on their religion, which is a clear indication of a hate speech. Such language can be extremely harmful and can lead to real-life violence and discrimination.



[Example 7]
[Instruction and Question]
In this task, you are given a public comment from online platforms. You are expected to classify the comment into two classes: identity-attack and non-identity-attack. Identity attack refers to anything rude or disrespectful about someone's identity, including appearance, race, gender, and sexual orientation.\n\nComment: Bail is a right under both Alaska and federal constitutions. A person charged with an offense is not legally guilty of it, no matter what the police say.  Curious, how some people burn incense before a constitutional right they like and ignore the ones they don't like.\n\nGranted, the court has discretion in setting the amount of bail, provided it is not \"excessive.\"  There are different interpretations of that word.  In some minds it means \"excessive compared to the bail usually set for similar charges.\"  In others, it means \"excessive compared to the defendant's resources.\"  Usually, the court considers both.\n\nPeople also don't realize that there is a distinction between cash bail and \"corporate\" bail.  The latter means a bail bondsman puts up a bond for you (for which he will mulct you dearly).  Usually, the defendant in felony cases is required to post at least 10% of his total bail amount in cash.  The article doesn't say whether the $2,500 was the whole bail, or only the cash part.

[Answer]
Non-identity-attack

[Rationale]
The comment does not contain anything rude or disrespectful about someone's identity, appearance, race, gender, or sexual orientation. So it could be classified into "Non-identity-attack".



[Example 8]
[Instruction and Question]
In this task you are given a tweet and you must identify whether the tweet contains any offense or any form of (untargeted) profanity. Label the post as NOT if the post does not contain offense or profanity.  Non-offensive posts do not include any form of offense or profanity. label the post as OFFENSIVE if the post contains offensive language or a targeted (veiled or direct) offense.  Twitter user mentions were substituted by @USER and URLs have been substitute by URL.\n\n@USER @USER For what? Why are you asking for his arrest?  Liberals being butthurt for 19 months is not an impeachable offense.

[Answer]
NOT

[Rationale]
The tweet does not contain any offense or profanity, but rather an opinion expressed in a civil manner. Therefore, the label for this tweet should be NOT.



[Example 9]
[Instruction and Question]
In this task, you are given a text from a social media post. Your task is to classify the given post into two categories: 1) yes if the given post is potentially offensive to anyone (i.e., a subset of people, any particular person, etc.), 2) no, otherwise. Note that potentially offensive posts can contain sexual, racial, religious biased or offensive language. Warning: the examples and instances may contain offensive language.\n\nKim K is alright but honestly so many niggas been in that I'd rather not.. And Miley got unfollowed that bitch is just strange.

[Answer]
Yes

[Rationale]
The given post is potentially offensive to some people as it contains a racial slur and objectifying language towards women. Therefore, the answer is "Yes".



[Example 10]
[Instruction and Question]

    \end{lstlisting}
    




\begin{lstlisting}[caption={Demonstration used for sni science}, label=sni science,,captionpos=b]
[Example 1]
[Instruction and Question]
In this task, you're given a paragraph from the research paper and your task is to generate a suitable title for the research paper based on the given paper. Under 100 words is a good title length.\n\nIn this study, we report the complete genome sequence of two contemporary human coronavirus OC43 (HCoV-OC43) strains detected in 2003 and 2004, respectively. Comparative genetic analyses of the circulating strains and the prototype HCoV-OC43 strain (ATCC VR759) were performed. Remarkably, a lower than expected similarity is found between the complete genomes and more in particular between the spike genes of the BE03 and BE04 strains. This finding suggests the existence of two genetically distinct HCoV-OC43 strains, circulating in Belgium in 2003 and 2004. Spike gene sequencing of three additional 2003 and two additional 2004 HCoV-OC43 strains, and subsequent phylogenetic analysis confirm this assumption. Compared to the ATCC prototype HCoV-OC43 strain, an important amino acid substitution is present in the potential cleavage site sequence of the spike protein of all contemporary strains, restoring the N-RRXRR-C motif, associated with increased spike protein cleavability in bovine coronaviruses. We here describe specific characteristics associated with circulating HCoV-OC43 strains, and we provide substantial evidence for the genetic variability of HCoV-OC43. D

[Answer]
Circulation of genetically distinct contemporary human coronavirus OC43 strains

[Rationale]
The title should provide evidence for the existence of two genetically distinct HCoV-OC43 strains, which were circulating in Belgium in 2003 and 2004, and describes specific characteristics associated with these strains. Also, one key aspect to consider in this paper is the genetic variability of HCoV-OC43 and the potential implications for increased spike protein cleavability. My suggestion would be: "Circulation of genetically distinct contemporary human coronavirus OC43 strains".



[Example 2]
[Instruction and Question]
In this task, you are given a sentence from the research paper and the category to which it belongs. Your task is to classify whether the given category is correct or not by providing \"True\" and \"False\", respectively. Here are the definitions for the categories: Background (Why is this problem important? What relevant works have been created before? What is still missing in the previous works? What are the high-level research questions? How might this help other research or researchers?), Purpose (What specific things do the researchers want to do? What specific knowledge do the researchers want to gain? What specific hypothesis do the researchers want to test?), Method (How did the researchers do the work or find what they sought? What are the procedures and steps of the research?), or Finding (What did the researchers find out? Did the proposed methods work? Did the thing behave as the researchers expected?). \",\n\nSentence: Systems biology offers considerable promise in uncovering novel pathways by which viruses and other microbial pathogens interact with host signaling and expression networks to mediate disease severity. \n Section: method

[Answer]
False

[Rationale]
The sentence describes the promise of systems biology in uncovering novel pathways by which viruses and other microbial pathogens interact with host signaling and expression networks to mediate disease severity. This is a background statement that highlights the importance of the problem and the relevance of the research area. Therefore, the correct category for this sentence is "Background", not "Method". Hence, the answer is False.



[Example 3]
[Instruction and Question]
Given an entity as input, output another entity which is part of the input entity. These are entities of meronym. In linguistics, meronymy is a semantic relation between a meronym denoting a part and a holonym denoting a whole. In simpler terms, a meronym (i.e., output entity) is in a part-of relationship with its holonym (i.e., input entity).\n\ngallinules

[Answer]
talon

[Rationale]
A gallinule is a type of bird, and one of its parts is the talon, which is a sharp claw on its foot used for grasping prey or perching on branches. Therefore, talon is a meronym of gallinules.



[Example 4]
[Instruction and Question]
Indicate the presence of numerical data in the abstract (whether numerical data plays a key role in medical records), solely using the information from the abstract.\n\nAbstract: A general approach to the child with foot pain has been described. In that the pediatrician is so often the recipient of parent's questions about children's shoes, he or she should be aware of problems relating to poorly fitting shoes, probably the most common source of complaints of foot pain in children. With a systematic approach, the pediatrician is very capable of managing the majority of such complaints. \nQuestion: Does the abstract contain any numerical data?

[Answer]
No

[Rationale]
It provides general information about foot pain in children and emphasizes the importance of proper shoe fit. However, it does not include any numerical values or measurements. So the answer is "No".



[Example 5]
[Instruction and Question]
In this task, you are given a paragraph, event and an entity. The event is part of the given paragraph and it changes the state of the entity. Your task is to classify the state of the entity into three classes: 1) not exist, 2) unknown location and 3) known location. \"not exist\" means the entity doesn't exist in that state anymore. \"unknown location\" means the entity exists in that state but location is unknown. \"known location\" means the entity exists and location is known. \n\n \"process paragraph : : Food enters the small intestine. The small intestine produces hormones. The hormones cause the pancreas to produce enzymes for digestion. The enzymes travel from the pancreas to the small intestine through tubes. The enzymes help the intestines digest the food. The pancreas produces the hormones insuline and glucagon. The hormones control the amount of sugar in the blood. \", \"event : The small intestine produces hormones.\", \"entity : food\"

[Answer]
known location

[Rationale]
The state of the entity "food" is still present and it is being processed by the small intestine and the enzymes produced by the pancreas, so its location is known to be in the small intestine. Therefore, the answer is "known location".



[Example 6]
[Instruction and Question]
In this task, you are given an abstract of article. Your task is to generate label \"True\" if abstract is structured, otherwise generate \"False\". A structured abstract is composed of a topic sentence (or key sentence), relevant supporting sentences, and a closing (or transition) sentence. This structure is key to keeping your abstract focused on the main idea and creating a clear and concise image.\n\nOBJECTIVE A20 is a TNF-inducible primary response gene, which has been found to have antiapoptotic function in several cancer cells. This study investigates A20 expression in human glioma tissues and four glioma cell lines, and its effect on tumorigenesis of glioma cells and a mouse tumor model.   \n METHODS Human glioma tissue samples and cells were subject to reverse transcription-PCR (RT-PCR), western blotting and immunohistochemistry. Glioma cells was tested by flow cytometry. A xenograft tumor model in mice was utilized to examine the knock-down effect of specific A20 siRNAs on tumorigenesis.   \n RESULTS A20 was overexpressed in clinical glioma tissue samples (63.9%) and correlated with clinical staging. All four human glioma cell lines expressed A20, among which U87 displayed the strongest expression signals. Inhibiting A20 expression by siRNAs in vitro reduced the growth rates of glioma cells and resulted in G1/S arrest and increased apoptosis. In a mouse tumor model, local administration of siRNA significantly suppressed solid tumor growth.   \n CONCLUSIONS A20 was overexpressed both in human glioma tissues and cell lines, and inhibiting A20 expression greatly slowed tumor cell growth in culture and in mice. These findings indicated that A20 is involved in tumorigenesis of human glioma, and may serve as a future therapeutic target.

[Answer]
True

[Rationale]
The abstract includes a topic sentence that states the objective of the study ("OBJECTIVE A20 is a TNF-inducible primary response gene, which has been found to have antiapoptotic function in several cancer cells. This study investigates A20 expression in human glioma tissues and four glioma cell lines, and its effect on tumorigenesis of glioma cells and a mouse tumor model."), relevant supporting sentences that describe the methods and results of the study, and a closing sentence that summarizes the main findings and conclusions of the study. This structure helps to keep the abstract focused and concise, and allows readers to easily understand the main idea and significance of the research. So the answer is “True”.



[Example 7]
[Instruction and Question]
In this task, you are given two statements. The task is to output whether a given textual premise, i.e. Statement 2, entails or implies a given scientific fact, i.e. Statement 1. The output should be 'entails' if Statement 2 supports Statement 1 and should be 'neutral' otherwise.\n\nSentence 1: Angiosperms is the term for flowering seed plants. Sentence 2: Angiosperms are also known as flowering plants because they produce flowers in which seeds are produced and in which they develop.

[Answer]
entails

[Rationale]
Statement 2 provides additional information that supports the scientific fact stated in Statement 1. Statement 1 defines angiosperms as "flowering seed plants", and Statement 2 explains that the reason they are called "flowering plants" is because they produce flowers that contain seeds and that serve as the site of seed development. Therefore, Statement 2 provides evidence that supports the definition of angiosperms provided in Statement 1. As a result, Statement 2 entails or implies Statement 1.



[Example 8]
[Instruction and Question]
You are given two sentences. You have to find if there is entailment or agreement of the Hypothesis by the Premise. From the given pair of sentences, you should identify if there is enough information in the Premise to support the claim made in the Hypothesis. The Premise may not exactly be the same as Hypothesis. Your task is to return 'entails' if the premise supports hypothesis else return 'neutral'.\n\nPremise: For example, airway resistance in the ivy leaf group decreased by 24% on day three of the study, compared to only 5% in the placebo group. \n Hypothesis: Poison ivy typically has three groups of leaves.

[Answer]
neutral

[Rationale]
The premise provides information about a study that measured airway resistance in the ivy leaf group and the placebo group. It does not provide any information about the number of groups that poison ivy leaves typically have. Therefore, there is no direct connection between the premise and the hypothesis, and the premise does not provide enough information to support or refute the claim made in the hypothesis. As a result, the output should be "neutral".



[Example 9]
[Instruction and Question]
In this task, you are given the abstract of a research paper. Your task is to generate a summary of this abstract. Your summary should not be very short, but it's better if it's not more than 30 words.\n\nThe description of neural computations in the field of neuroscience relies on two competing views: (i) a classical single-cell view that relates the activity of individual neurons to sensory or behavioural variables, and focuses on how different cell classes map onto computations; (ii) a more recent population view that instead characterises computations in terms of collective neural trajectories, and focuses on the dimensionality of these trajectories as animals perform tasks. How the two key concepts of cell classes and low-dimensional trajectories interact to shape neural computations is however currently not understood. Here we address this question by combining machine-learning tools for training RNNs with reverse-engineering and theoretical analyses of network dynamics. We introduce a novel class of theoretically tractable recurrent networks: low-rank, mixture of Gaussian RNNs. In these networks, the rank of the connectivity controls the dimensionality of the dynamics, while the number of components in the Gaussian mixture corresponds to the number of cell classes. Using back-propagation, we determine the minimum rank and number of cell classes needed to implement neuroscience tasks of increasing complexity. We then exploit mean-field theory to reverse-engineer the obtained solutions and identify the respective roles of dimensionality and cell classes. We show that the rank determines the phase-space available for dynamics that implement input-output mappings, while having multiple cell classes allows networks to flexibly switch between different types of dynamics in the available phase-space. Our results have implications for the analysis of neuroscience experiments and the development of explainable AI.

[Answer]
A theoretical analysis of a new class of RNNs, trained on neuroscience tasks, allows us to identify the role of dynamical dimensionality and cell classes in neural computations

[Rationale]
Here's a brief summary:\nA theoretical analysis of a new class of RNNs, trained on neuroscience tasks, allows us to identify the role of dynamical dimensionality and cell classes in neural computations.



[Example 10]
[Instruction and Question]

    \end{lstlisting}
    




\begin{lstlisting}[caption={Demonstration used for sni mcqa}, label=sni mcqa,,captionpos=b]
[Example 1]
[Instruction and Question]
We would like you to assess the QUALITY of each of the following argument (discussing Gay Marriage) and determine if the argument is Valid or Invalid. A valid argument is clearly interpretable and either expresses an argument, or a premise or a conclusion that can be used in an argument for the topic of gay marriage. An invalid argument is a phrase that cannot be interpreted as an argument or not on the topic of gay marriage.\n\nSince marriage (love) is legally sanctioned for all heterosexuals, the same rights and protection should be granted to homosexuals who choose to marry...or these \"legalities\" should be terminated completely.

[Answer]
Valid

[Rationale]
The argument presents a clear argument for granting the same legal rights and protections to homosexuals who choose to marry as those granted to heterosexuals who marry, or for terminating legal marriage altogether.



[Example 2]
[Instruction and Question]
In this task, you will be shown an English sentence. You need to classify the sentence as either a representation of an anaphor number agreement or as an incorrect representation. An anaphor is an expression whose interpretation depends upon another expression. Anaphor number agreement is a restriction where a phrase agrees with the preceeding expression: plurals match plurals or singular expressions match singular expressions. In this task a valid representation of anaphor number agreement should be denoted as 'good' while everything else should be denoted as 'bad'. A correct representation of an anaphor number agreement will be a grammatically correct instance of a sentence where the subject and self-referencing nouns match plurality. An incorrect representation is a sentence in which the subject and self-referencing noun's plurality do not match.\n\nMost children didn't boast about herself.

[Answer]
bad

[Rationale]
The subject of the sentence is "Most children", which is plural. However, the reflexive pronoun "herself" is singular and does not agree with the plurality of the subject. A correct representation would be "Most children didn't boast about themselves", where the plural subject "children" agrees with the plural reflexive pronoun "themselves".\nTherefore, This sentence is an incorrect representation of anaphor number agreement and should be denoted as "bad".



[Example 3]
[Instruction and Question]
In this task, you are given a short passage that conveys a stereotype. A stereotype is an over-generalized belief about a particular group of people. You are expected to classify the passage into nine types of stereotype: race-color, socioeconomic, gender, disability, nationality, sexual-orientation, physical-appearance, religion, and age.\n\nPassage: He always took care of the yard work himself. He operated a riding lawn mower like no one's business.

[Answer]
gender

[Rationale]
The use of the pronoun "he" and the description of the subject's skill with the lawn mower are indicative of this gender stereotype. Therefore, the passage can be classified as a gender stereotype.



[Example 4]
[Instruction and Question]
Classify the given a piece of financial news into three classes: positive, negative, and neutral. Output must be 'positive', 'negative', or 'neutral'.\n\nFinnish department store chain Stockmann Oyj Abp net profit rose to 39.8 mln euro ( $ 56.8 mln ) for the first nine months of 2007 from 37.4 mln euro ( $ 53.4 mln ) for the same period of 2006 .

[Answer]
positive

[Rationale]
This financial news states that the net profit of Finnish department store chain Stockmann Oyj Abp has increased from 37.4 million euros to 39.8 million euros for the first nine months of 2007 as compared to the same period of 2006. This indicates a growth in profits for the company, which is generally considered as a positive development. So, this financial news can be clarified as "positive".



[Example 5]
[Instruction and Question]
In this task, you will be shown a prompt from a judicial decision and multiple holding statements derived from citations following text in a legal decision.  Holdings represent the governing legal rule when the law is applied to a particular set of facts. There are five answer choices for each citing text. The correct answer is the holding statement that corresponds to the citing text. The four incorrect answers are other holding statements. You should find the correct option. There is a <HOLDING> token in the position of the citing text prompt where the holding statement was extracted.\n\npersons for damages that result because of an injury to the corporation unless a special injury exists in the form of either: (1) a special duty, such as a contractual duty, between the defendant and the shareholder; or (2) an injury separate and distinct from that suffered by other shareholders. Healthsource, Inc., 2005-NMCA-097, \u00b6 25, 138 N.M. 70, 116 P.3d 861; Marchman, 120 N.M. at 81-82, 898 P.2d at 716-17. {14} Defendants argue that Marchman and Healthsource, Inc. are controlling and that Plaintiff lacks standing because neither exception applies. Furthermore, Defendants argue that, to the extent Plaintiff alleges mismanagement or negligence, these claims may also be asserted only derivatively. See Schwartzman v. Schwartzman Packing Co., 99 N.M. 436, 441, 659 P.2d 888, 893 (1983) (<HOLDING>). The district court agreed, and dismissed \n Holding statements: (A) recognizing privilege for corporate officers directors and shareholders to influence the actions of their corporation (B) holding that claims of corporation vest in corporation (C) holding that the trust fund and denuding theories of recovery were considered to belong to the debtor corporation because each action was created for the benefit of the corporation ie to vindicate injury to the corporation caused by improper actions by control persons (D) holding that a corporate plaintiffs claims must be dismissed on the ground of failure of any attorney to appear for the corporation (E) holding that claims against corporate officers for mismanagement belong to the corporation

[Answer]
(E)

[Rationale]
The defendants argue that the plaintiff lacks standing and claims of mismanagement can only be asserted derivatively. The holding statement cited in support of this argument is "holding that claims against corporate officers for mismanagement belong to the corporation" which is option (E).



[Example 6]
[Instruction and Question]
In this task, you are given a multiple-choice question about healthcare. Answer the question based on your information and classify your answers into '1', '2', '3', and '4'.\n\nQuestion: The chemical reactions of the splicing process of the introns during the maturation of the pre-mRNAs consists of: \n Options:  <1> An oxidation-reduction reaction. <2> A transesterification reaction. <3> Two sequential transesterification reactions. <4> Three sequential transesterification reactions. <5> A transesterification reaction followed by another oxidation reaction.

[Answer]
3

[Rationale]
The splicing process of the introns during the maturation of pre-mRNAs involves the removal of introns and joining of exons to form a mature mRNA molecule. This process is catalyzed by a large ribonucleoprotein complex called the spliceosome, which carries out two sequential transesterification reactions to cleave the intron-exon junctions and ligate the exons together.



[Example 7]
[Instruction and Question]
Given a passage classify if the passage has a definite objective/aim/goal or not. Output '1' if the passage has a defininte objective/aim/goal and output '0' if the passage does not have a definite objective/aim/goal.\n\nThe majority of studies assessing executive function in attention deficit disorder (ADD) have shown deficits in attentional set shifting using either the Wisconsin card sorting task or the intra-dimensional/extra-dimensional set-shifting task (ID/ED). Damage to the prefrontal cortex in humans, primates, and rodents impairs extra-dimensional (ED) shifts. Noradrenergic depletion of the medial prefrontal cortex in rats is sufficient to impair attentional set shifting. Atomoxetine, a selective norepinephrine (NE) re-uptake inhibitor, is hypothesized to produce beneficial effects in patient with ADD by augmenting NE release in prefrontal cortex.', 'We assessed the effects of systemic administration of atomoxetine (0.0, 0.1, 0.3, and 0.9 mg/kg/ml) in normal and noradrenergically lesioned (NE-LX) rats on attentional-set shifts. We replicated findings showing NE-LX rats are selectively impaired on the ED shifts but not reversals or other discriminations.', 'Atomoxetine remediated the attentional set-shifting impairments in NE-LX rats but impaired ED performance of non-lesioned rats.

[Answer]
0

[Rationale]
The passage describes a series of studies that investigate the effects of various manipulations on attentional set-shifting in rats, but it does not state a specific research question or hypothesis that the studies are designed to test. The passage provides information on the methods, results, and interpretations of the studies, but it does not explicitly state a main objective or aim. Therefore, the output is '0'.



[Example 8]
[Instruction and Question]
You will be given a topic and an argument. Decide the argument's stance towards that topic. The argument's stance is in favor or against the topic. If the argument supports that topic, answer with \"in favor\"; otherwise, if the argument opposes the topic, answer with \"against\".\n\ntopic: Tibet independence\nargument: Tibet can follow the East Timor model to achieving independence

[Answer]
in favor

[Rationale]
The passage suggests that Tibet can achieve independence by following the model used by East Timor. This implies that the author believes that Tibet should be independent and supports the idea of Tibet becoming an independent nation. Therefore, the argument is "in favor" of Tibet independence.



[Example 9]
[Instruction and Question]
In this task, you are given a word, followed by two sentences. Your task is to figure out whether both the sentences use the aforementioned word with the same meaning. You should respond with 'True' if the words in both sentences share the same meaning, and 'False' otherwise.\n\ncommunicate \n Sentence1: Many deaf people communicate with sign language. \n Sentence2: I feel I hardly know him; I just wish he'd communicate with me a little more.

[Answer]
False

[Rationale]
The word "communicate" is used with different meanings in both sentences. In the first sentence, "communicate" refers to the act of conveying information or expressing oneself through sign language, while in the second sentence, it refers to the act of sharing thoughts or feelings with someone else. The two meanings are different, so the answer is False.



[Example 10]
[Instruction and Question]

    \end{lstlisting}
    




\begin{lstlisting}[caption={Demonstration used for sni math}, label=sni math,,captionpos=b]
[Example 1]
[Instruction and Question]
You are provided with an arithmetic question. Your task is to compute the solution using the given arithmetic operations. The only arithmetic operators needed to answer the questions are'+'(addition) and'-'(subtraction). The answer should be correct to one decimal place.\n\nKeith spent $6.5 on a rabbit toy, $5.79 on pet food, and a cagecost him $12.51. He found a dollar bill on the ground. What was the total cost of Keith's purchases?

[Answer]
24.8

[Rationale]
Keith spent $6.5 on a rabbit toy, $5.79 on pet food, and a cagecost him $12.51. Finding a dollar doesn't matter with Keith's purchase. (Total cost of Keith's purchase) = 6.5 + 5.79 + 12.51 = 24.80.



[Example 2]
[Instruction and Question]
Given a math problem with context and a question and 5 answer choices, the task is to provide the correct answer choice based on the problem. You must choose one of the given answer choices by letter: a, b, c, d, or e; anything else is invalid.\n\nProblem: 20 percent of country y ' s yearly exports come from fruit exports. one - sixth of all fruit exports from country y are orange exports. if country y generates $ 4.25 million from its orange exports, how much money does it generate from its yearly total of exports ?\nOptions: a. $ 21.25 m, b. $ 25.5 m, c. $ 106.25 m, d. $ 127.5 m, e. $ 153 m

[Answer]
d

[Rationale]
To solve the problem, we can use the following steps:\n1. Find the total revenue from fruit exports by dividing the revenue from orange exports by the fraction of orange exports out of all fruit exports:\n4.25 million / (1/6) = 25.5 million\n\n2. Use the fact that fruit exports make up 20% of the total exports to find the total revenue from all exports:\nTotal exports = (25.5 million) / (0.2) = 127.5 million



[Example 3]
[Instruction and Question]
Given a simple high-school level math question, you are required to solve it and provide the final answer. The final answer is always a single number. These questions can range from a variety of topics like simple arithmetic, solving equations, converting a quantity from one unit to another, finding remainders/GCD/LCM, finding probabilities etc. Each question has only one correct answer. This answer can be a positive or negative integer, a fraction or a decimal number. If the answer is a negative number use the hyphen (e.g. -42) symbol for the minus sign. For decimal numbers, do not add extra zeros after the decimal point. For fractional numbers, separate the numerator and denominator using a forward slash (e.g. 3/25).\n\nThree letters picked without replacement from svojjvzsoovvso. Give prob of sequence jvs.

[Answer]
1/91

[Rationale]
The total number of ways to pick 3 letters from the given string is given by the combination formula, which is C(14,3) = 364. There are 2 j's, 3 v's ,2 s's in the sequence. So the number of possible cases to make sequence "jvs" would be 2x3x2=12. The final answer is 12/364 which can be simplified as 1/91.



[Example 4]
[Instruction and Question]
In this task, you need to answer the given multiple-choice question on the general math. Classify your answers into 'a', 'b', 'c', 'd', and 'e'.\n\nProblem: in track last week , the boys ran 34 laps . the girls ran 20 more laps . each lap is a 1 / 6 th of a mile . how many miles did the girls run ? \nOptions: a ) 9 miles , b ) 10 miles , c ) 11 miles , d ) 12 miles , e ) 13 miles

[Answer]
a

[Rationale]
The girls ran 34 + 20 = 54 laps in total. Each lap is 1/6 of a mile, so the girls ran 54 x 1/6 = 9 miles.



[Example 5]
[Instruction and Question]
In this task, you are given a statement spoken by a politician in natural language and the subject of the discussion for the given statement. Your task is to determine if the given subject of the discussion is correct for the given statement. The subject can be multiple words. Answer with \"True\" if it's the correct subject of discussion or return \"False\" if it's incorrect.\n\nStatement: \"My plan offers three times as much tax relief to the middle class as Senator McCain's.\" \n Subject: taxes

[Answer]
True

[Rationale]
The subject of discussion in the statement is taxes, and the politician is comparing their plan's tax relief to Senator McCain's. Therefore, the answer is True.



[Example 6]
[Instruction and Question]
You are given a mathematical question described with a circumscribed mathematical vocabulary. Questions in this task usually contain some variables (e.g. X, Y), equations (e.g. X+Y=3) and some relationships (X is greater than Y). Often you need to find the value of the variable using equations and relationships. You are also given 4 or 5 answer options (associated with \"A\", \"B\", \"C\", \"D\", \"E\"). Do not generate anything else apart from one of the following characters: 'A', 'B, 'C', 'D', 'E'. LaTeX mathematical format (the standard way to express mathematical expressions in the typesetting software known as LaTeX) is used to express equations. Each question is solvable with high school math knowledge. Give only one answer for each question.\n\nFor the function \\(f\\), \\(f(x)\\) is inversely proportional to \\(x\\). If \\(f(k) = a\\), what is the value of \\(f(\\frac { k } { 3 } )\\) ?\n(A)\\(\\frac { a } { 3 } \\) (B)\\(\\frac { 2a } { 3 } \\) (C)\\(3a\\) (D)\\(6a\\) (E)\\(9a\\)

[Answer]
C

[Rationale]
If \(f(x)\) is inversely proportional to \(x\), then \(f(x)\cdot x = k\) for some constant \(k\). If \(f(k) = a\), then we have \(f(k)\cdot k = a\cdot k), so \(a\cdot k = k\cdot f(k) = k\cdot f(\frac{k}{3})\cdot 3 = 3f(\frac{k}{3})\cdot k/3\). Simplifying this gives us \(a = f(\frac{k}{3})\cdot k\), so \(f(\frac{k}{3}) = a/k\cdot 3 = 3a/k\). Therefore, the answer is (C) \(3a\).



[Example 7]
[Instruction and Question]
You are given a math word problem and you are supposed to apply an addition mathematical operator on the numbers embedded in the text to answer the following question and then only report the final numerical answer.\n\nContext: Edward spent $ 9. Then he spent $ 8 more. Now he has $ 17. \n\n Question: How much did Edward have before he spent his money?

[Answer]
34

[Rationale]
To solve the problem, we need to use subtraction, not addition.\nSince Edward spent $9 and then spent an additional $8 more, his total spending was $9 + $8 = $17.\nWe know that Edward currently has $17, so to find out how much money he had before he spent it all, we need to add $17 with his total spent money $17.




[Example 8]
[Instruction and Question]
In this task you will be given an arithmetic operation in Italian and you have to find its answer. The operations 'addition' and 'subtraction' have been replaced with their italian translations i.e you need to perform addition when you see 'aggiunta' and subtraction in case of 'sottrazione'.\n\n1195 sottrazione 570 aggiunta 7850 aggiunta 8195 sottrazione 2484

[Answer]
14186

[Rationale]
1195 - 570 + 7850 + 8195 - 2484 = 625 + 7850 + 8195 - 2484 = 8475 + 8195 - 2484 = 16670 - 2484 = 14186.



[Example 9]
[Instruction and Question]
A ploynomial equation is a sum of terms. Here each term is either a constant number, or consists of the variable x raised to a certain power and multiplied by a number. These numbers are called weights. For example, in the polynomial: 2x^2+3x+4, the weights are: 2,3,4. You can present a polynomial with the list of its weights, for example, equation weights = [6, 4] represent the equation 6x + 4 and equation weights = [1, 3, 4] represent the equation 1x^2 + 3x + 4. In this task, you need to compute the result of a polynomial expression by substituing a given value of x in the given polynomial equation. Equation weights are given as a list.\n\nx = 3, equation weights = [6, 3, 2]

[Answer]
65

[Rationale]
The given polynomial equation is: 6x^2 + 3x + 2\nSubstituting x = 3, we get: 6(3)^2 + 3(3) + 2 = 54 + 9 + 2 = 65



[Example 10]
[Instruction and Question]

    \end{lstlisting}
    




\begin{lstlisting}[caption={Demonstration used for sni exqa}, label=sni exqa,,captionpos=b]
[Example 1]
[Instruction and Question]
Given a paragraph about movies, and a set of conversational question answers about the paragraph, generate a relevant question to the topic of the paragraph. The paragraph has the prefix 'CONTEXT:'. Each conversation question has a prefix `Q:` followed by the answer prefix `A:`.\n\nCONTEXT: They didn't change the whole climax scene. In one version Rudra's father saying him that the girl he loved is actually his step sister, so he can not marry her. On the other hand, the scene is same but there is no dialogue as it is muted and audience needs to understand from their reaction and background score. This change was made because the film got poor audience response due to the dialogue delivery of the conversation.  <sep>  Q: What are the differences between two climax of Solo? A: In one version Rudra's father saying him that the girl he loved is actually his step sister, so he can not marry her

[Answer]
What happens in the other version?

[Rationale]
The given conversation question provides information about one version of the climax scene in the movie, but does not provide any information about the other version. Therefore, asking about what happens in the other version can provide more information about the differences between the two versions of the climax scene.



[Example 2]
[Instruction and Question]
In this task, you're given a question, along with a context passage which has extra information available on certain terms mentioned in it, i.e., the proper nouns in the passage. Your job is to determine whether information from more than one term is needed to answer the question. Indicate your choice as `a` for Yes or `b` for No.\n\nQuestion: Who was the designer of the Swordfish? Passage:As designed in 1931, the He 46 was a two-seat sesquiplane of mixed construction. The upper wing was swept back 10\u00b0. This is a common trick when a change in an aircraft's design moves the center of gravity to the rear; rather than moving the whole wing back to counter for the shift, it is easier to simply angle the wings back slightly, which requires very little change to the basic design, or even the overall structure of the wing. (The similarly angled wings on the Fairey Swordfish are due to just such a mid-program change to the weight distribution of the aircraft.) The tailplane was mounted high and braced by struts. The undercarriage was fixed, and the tail was fitted with a skid rather than a wheel. The He 46 prototype first flew in late 1931; its flight characteristics were good, but design improvements were incorporated. The small lower wing was removed, while the mainplane was increased in area by 22% and braced to the fuselage, transforming the He 46 into a parasol-wing monoplane. A more powerful engine was added to the second prototype, and a single 7.92\u00a0mm (.312\u00a0in) MG 15 machine gun for the rear seat was added to the third prototype.\n

[Answer]
b

[Rationale]
The question only asks for the designer of the Swordfish, which is mentioned directly in the passage as "The similarly angled wings on the Fairey Swordfish." No other information or proper nouns are needed to answer the question.



[Example 3]
[Instruction and Question]
Indicate a text span in the passage about which an inquisitive question is asked. (note, don't answer the question).\n\nSentence: In trading on the Stock Exchange of Hong Kong , the shares closed Wednesday at HK $ 4 . 80 each .<sep>Question: Is there any fee associated with owning foreign shares?

[Answer]
Stock Exchange of Hong Kong

[Rationale]
The text span should indicate a specific location where the shares were traded and the question pertains to the fee associated with owning shares in foreign markets. The relevant text span is "Stock Exchange of Hong Kong".



[Example 4]
[Instruction and Question]
Given a passage and a query based on the passage, generate an unambiguous, concise and simple answer to the query from information in the passage. Note that the answer may not be present in exact form.\n\nBuy chicken before this date for storage in your home refrigerator. Always cook or freeze the chicken within two days of bringing it home, regardless of what the package date tells you. Never store raw chicken in the refrigerator at home for longer than two days. Once you've cleared up what belongs where, the next step is knowing how long the food will last. Storing food is a delicate matter -- in the fridge, food might get moist and moldy, and in the freezer it could lose all its moisture or get freezer burn. If you leave food in the fridge or freezer too long, it could not only lose its integrity but it may also become unsafe to eat. FoodSafety.gov has two useful charts listing how long certain foods will stay good in both the refrigerator and the freezer. The first chart lists all kinds of commonly refrigerated foods and the second is specifically for eggs and food made with eggs. Chrissy Metz, star of <strong>This is Us</em>, appeared on <strong>Jimmy Kimmel Liv... 1  Chrissy Metz, star of This is Us, appeared on Jimmy Kimmel Live on Tuesday wearing a stylishly on-point outfit that people are dying to copy\u2014and it\u2019s possible to do so without breaking the bank. References. 1  Basics for handling food safely. U.S. Department of Agriculture. 2  Foodborne illness: What consumers need to know. U.S. Department of Agriculture. 3  Food safety counts! U.S. Department of Agriculture. 4  Dietary Guidelines for Americans, 2010. U.S. Department of Health and Human Services. 5  Slow cookers and food safety. Keep your refrigerator below 40 degrees Fahrenheit, and store raw chicken in the coldest part of the fridge, which is usually the back of the bottom shelf or in the meat drawer. Make sure to decrease the storage time if you can't guarantee that the temperature will stay under 40 degrees F. (Also) I feel like I\u2019m limited to what I can do with the chicken once its cooked. Thanks for your help, SS. Dear SS: According to Foodsafety.gov, whole, uncooked poultry (chicken or turkey) is safe in the refrigerator for 1-2 days and safe in the freezer for up to 12 months. Poultry pieces are safe in the refrigerator for 1-2 days and safe in the freezer for up to nine months. Cooked poultry is safe refrigerated for 3-4 days and in the freezer for 2-6 months. You can also freeze your cooked breaded chicken cutlets (Foodsafety.gov specifically lists nuggets and patties) for 1-3 months. How long can you keep chicken in the fridge? Raw chicken: You'll need to either cook or freeze it fast. According to recommendations from the U.S. Food and Drug Administration, raw chicken will only keep in the refrigerator for about 1-2 days. (Same goes for turkey and other poultry.) If you store it in the freezer in a sealed, freezer-safe container or vacuum-sealed package, it should stay safe to thaw and eat for up to 9 months, in pieces, or up to 12 months, for a whole chicken. Cooked chicken: You have a little more time, but not much. I have to say that I have kept chicken in the fridge for longer than 2 days and not died. Some even have best before dates 3-4 days in. The quality definitely suffers, and I'm sure there's some risk, but it's probably low. How long does raw chicken last after being frozen and thawed? Chicken that has been defrosted in the fridge can be kept for an additional 1 to 2 days in the refrigerator before cooking; chicken that was thawed in the microwave or in cold water should be cooked immediately. I didn't even realize it until the next day when I found the container with the fresh chicken and realized that there were two containers of it in the fridge because I'd forgotten to throw away the old chicken from the previous week. Query: how long is chicken good in the refrigerator

[Answer]
According to Foodsafety.gov, whole, uncooked poultry (chicken or turkey) is safe in the refrigerator for 1-2 days

[Rationale]
Here's an direct answer: According to Foodsafety.gov, whole, uncooked poultry (chicken or turkey) is safe in the refrigerator for 1-2 days.



[Example 5]
[Instruction and Question]
You are given a sentence and a question, construct 2 answer options in a specific format i.e. ['option1','option2']. The answer options have to be text spans from the given sentence. Only one among the options should be the correct answer to the given question. Note that, the correct answer to the given question will require understanding of coreference resolution. Coreference resolution is the task of clustering mentions in text that refer to the same underlying real world entities.\n\nSentence: In July, Kamtchatka declared war on Yakutsk. Since Yakutsk's army was much better equipped and ten times larger, they were victorious within weeks. Question: Who was victorious

[Answer]
['Yakutsk', 'Kamchatka']

[Rationale]
The correct answer to the given question is 'Yakutsk', as they were the ones who won the war. However, to understand this, coreference resolution is required, as 'they' refers to Yakutsk in the sentence. So, the answer options should be: ['Yakutsk', 'Kamtchatka'] where 'Yakutsk' is the correct answer and 'Kamtchatka' is the incorrect option.



[Example 6]
[Instruction and Question]
Given a text paragraph and a fact, generate 'True' if the fact can be inferred/concluded from the paragraph, otherwise generate 'False'. The input passage contains as set of logical statements so the fact can be marked either 'True' or 'False' if the relavant info regarding the fact is present in the given passage. If the Fact can't be marked True or False from the given paragraph, mark it as False. Don't use knowledge outside the given paragraph to check the fact.\n\nParagraph: If something is rough then it is big. Gary is kind. Erin is rough. Erin is red. Gary is green. Gary is blue. Erin is blue. If Gary is rough then Gary is green. All kind things are blue. All quiet things are big. Bob is blue. Bob is rough. If something is green and blue then it is quiet. Bob is green. Red, quiet things are blue. Bob is big.     Question: Erin is not red.

[Answer]
False

[Rationale]
There is a statement in the paragraph that indicates Erin is red.  Therefore, the fact that "Erin is not red" can't be inferred from the paragraph.



[Example 7]
[Instruction and Question]

    \end{lstlisting}
    




\begin{lstlisting}[caption={Demonstration used for sni ethics}, label=sni ethics,,captionpos=b]
[Example 1]
[Instruction and Question]
In this task, you are given two simple actions (associated with \"A\", \"B\"). You must identify which action is considered less ethical. Do not generate anything else apart from one of the following characters: 'A', 'B'.\n\n(A) telling my friend to stop venting about her breakup with me\n(B) telling my bf that in case I get pregnant, I get to make the last call on whether we are keeping it or not

[Answer]
A

[Rationale]
For action A, telling a friend to stop venting about their breakup may be seen as unsupportive and insensitive, but it does not involve making a significant decision that affects another person's life. It is a more minor offense.\nFor action B, telling a partner that you will make the final decision about whether to keep a potential child or not, without their input or agreement, is a significant violation of their autonomy and a potentially manipulative or coercive action.\nTherefore, action B has a greater ethical implication and is considered less ethical than action A.



[Example 2]
[Instruction and Question]
In this task you are given an anecdote, which recount a complex ethical situation. The question you are given is, has the author of the anecdote acted ethically? answer this question with \"RIGHT\" or \"WRONG\".\n\nTLDR at the bottom.\n\n&amp;#x200B;\n\nThis did happen a few years ago but to this day it still haunts me. I met a girl at work (lets call her T), and she was perfect. We went out a few times and everything was good. Unfortunately T had to move away for school, not super far but far enough that we couldn't hang out to much. \n\nNow T and her sister (we can call the sister K) both worked with me. I was not the biggest fan of K but she was the sister of T, and i was determined to be in good graces with her, as Those two were very close.  So, maybe a week after T went back to school, I was working a shift with K and another friend (we can call him Q) the first half of the shift was fine, people pouring in and we barely had time to talk. Things took a turn when we all went on break together. I was craving a burrito and there was a Starbucks near by where I was going, so I offered K a ride since she was going there. We get back and Q is sitting there in the break room. \n\nWe sit down and Q looks at me and asks if I want to see this YouTube video of \"The Ultimate Vape\". I'm bored and have nothing better to do so I agree. We get to the point where its time to choose the flavor for the vape, and its churro flavor. This is what sets off the argument, because the man in the video didn't roll the R in churro (this is important to K because her and T are both half Black half Mexican women).  this is the conversation that followed.\n\nQ: Some people just can't roll their R's.\n\nK: That's why i hate White people,\n\nQ: ... that's kind of racist...\n\nThat's all i can fully remember was said but the events that followed were K saying you can't be racist towards White people and that's what THEIR argument was about. I considered both K and Q to be my friends but they were not necessarily friends with each other. So I sit aside let them have at it because if I choose a side I'm definitely loosing one of them as a friend.\n\nAt some point K turns to me, and she says,\n\nK: hey LexeL... I forgot my phone in my car, can you use yours to look up the definition of Racism?\n\nMe: I'm sorry K but this is between you two and id rather not get in the middle.\n\nAnd then she goes back to Q. i should mention that arguing over the Dictionary definition of racism had gone on for 3 minutes prior to this conversation. and Q had correctly said, Word for Word, what the definition of Racism was but K would not accept it.\n\nWe can skip to the end of break. I'm getting ready to go downstairs and K stops me.\n\nK: Have you talked to T lately?\n\nMe: Well I messaged her on Friday, but I know shes busy so I don't expect to talk to her until we Skype tomorrow. Why?\n\nK: Oh... well i just would expect to hear from her at all.\n\nThis interaction led me to believe that K had talked to T way before her argument with Q, and that she knew something I didn't, I spent the rest of my shift, nerves on end hoping I'd get to the bottom of this. \n\nK got off work before I did, terrible start to the rest of my night TBH. So as soon as I get off work this is the following conversation that I know by heart.\n\nText messages -\n\nT: what are you playing at??\n\nMe: Listen, idk what you've heard, who you've talked to, or what you're thinking, but I think we should talk.\n\nT: I agree.\n\nMe: Are you able to Skype tonight?\n\nT: We should just call on the phone...\n\nMe: okay.\n\nI have quite the drive home and I call her as soon as its over, we have a long talk about the situation that happened and she says that my story lines up with K's story. But... sadly I was in the wrong. Because I took no side in the situation, I actually took Q's side. \"a bystander is just taking the oppressors side without vocalizing it\" T said that she didn't think we should continue talking. We said goodbye for the last time and I never really got over it. so Am i the Asshole for not choosing her sisters side when her sister was wrong?\n\nTLDR; The sister of the girl I was seeing and another friend got into an argument about the definition of racism during break and i didn't choose a side, thus making me the bad guy in the eyes of the girl i was seeing.\n\nI wanna thank anyone for reading, even the TLDR, and I appreciate any comments. I'm just looking to learn from a past experience.\n\nThanks again!

[Answer]
RIGHT

[Rationale]
The author chose to remain neutral and not take a side in the argument, which led to the girlfriend's sister feeling unsupported and angry with the author. While it is understandable that the girlfriend's sister felt hurt and unsupported, the author acted ethically by not taking sides in the argument. Choosing a side in an argument between two friends can be tricky and can lead to further conflicts and hurt feelings. By choosing not to take a side, the author avoided causing more harm to either friend, and instead encouraged them to work out their differences and come to a resolution on their own.\nThe author's decision to remain neutral was a responsible and ethical choice that prioritized the well-being of both friends. Therefore, the answer to the question ""Has the author acted ethically?"" is ""RIGHT"".



[Example 3]
[Instruction and Question]
In this task you are given an anecdote. You must find out whether the anecdote is something the author has done or is considering doing. Label the instances as \"HISTORICAL\" when the author has already done something and label them as \"HYPOTHETICAL\" when the author is considering doing something in the future.\n\nSo I dated a girl for a few months, it was shaping up to be a longterm thing, but things went bad. We basically went a few months with no contact, until she contacted me, saying she was a few months pregnant and needed my support. I was floored. We'd barely spoken since the breakup, and I was clear as day about the fact that I NEVER wanted kids. \n\nThis was one of the worst things that could conceivably happen to me. She suggested getting back together, which I shut down right out of the gate. I explained that we were through months ago, and that I don't want to be a father, and hung up. She is dead set on keeping it, there have been tests, and the kid is mine, though I dont know how the fuck it happened, since we were fairly strict with contraceptives.\n\nI have a new GF, and we've just moved to a great new apartment. This phonecall literally shook me to my core, and I am freaking out like crazy. I've been searching for every possible way out, and I heard of a financial abortion, or whatever its called. Basically I want to resign any rights and responsibilities for the kid, since I never wanted him/her to begin with. I refuse to pay out the ass for a kid I don't want, and I obviously won't and can't force the mother to get rid of the baby, or give it up for adoption.\n\nHer friends have called me, saying that I really hurt her by refusing to acknowledge this kid, and calling me selfish and whatnot, but I can't handle all that stuff right now, or ever probably. Am I the asshole for wanting a \"financial abortion\", that is, no contact, no child support, no custody, and is it morally and legally okay???

[Answer]
HISTORICAL

[Rationale]
The author describes a past event in which they dated someone, broke up, and then received news that their ex-girlfriend was pregnant. The author expresses reluctance to be a father and considers ways to resign from any rights and responsibilities for the child, which suggests that the event has already occurred. Therefore, the label for this anecdote is "HISTORICAL".



[Example 4]
[Instruction and Question]
Given a post that is a real-life anecdote of a complex ethical situation and an associated claim about its type, verify if the claim is true or not. The claim asks if the posts are historical or hypothetical. The posts are \"HISTORICAL\" when the author has already done something and they are \"HYPOTHETICAL\" when the author is considering doing something in the future. Say 'yes' if the claim is true otherwise say 'no'\n\nCLAIM : The following post is HISTORICAL. POST : Please excuse spelling I'm on mobile. So first off I would like to let everyone know that I am in a theater  program and it was two days before the show starts (it starts on the 29th) and we had beed rehearsing since after lunch until 5:30. Some kids were whispering during the rehearsal and some others were on their phone backstage. I wasn't on my phone or talking unless I was telling someone else to be quiet.\n     Our theater teacher (who we will call R) was angry at people for not paying attention. R might have just been stressed, but so were we and being yelled at after finishing a two hour show rehearsal was not what we needed. My friend, who we will call J, was there too and after he had finished yelling at everyone I asked her, and I quote, \"does he always get like this before a show?\", and she called me an asshole. Please don't say anything mean about J because she is still my friend, and R was just stressed. So Reddit AITA?\n\nTLDR: My theater teacher yelled at us two days before a show and I asked my friend if he always gets like this before a show, and my friend called me an asshole.

[Answer]
no

[Rationale]
The author of the post is questioning if it was appropriate to ask their friend if their theater teacher always behaves in such a manner when under stress.\nIn the post, the author uses the past tense to describe the events, but they do not explicitly state when the events occurred. However, they mention that the show starts on the 29th, which suggests that the events occurred in the past, before that date.\nBased on the above analysis, it seems that the post is not hypothetical, as the author describes events that have already occurred. Therefore, the appropriate category is ""HISTORICAL."" In that case, the answer would be ""no"".



[Example 5]
[Instruction and Question]
Given a real-life anecdote of a complex ethical situation, judge whether or not a specific person is wrong here. Answer 'not answerable' if more information is required and 'answerable' if all the information is sufficient to answer who is wrong in this situation.\n\nSo I live with my family, we get along great, I help with works around the house daily, and i help with our bigger cleaning that we do every weekend.\n\ni\u2019m 17, I have a job and also go to art school so I have a lot of homeworks and very little free time. I have a boyfriend and a group of friends.\n\nI usually spend my fridays at the bf and then go home.\n\nMy mom seems to have a problem with this, she always says how I\u2019m a bad daughter because I\u2019d rather be with my friends/bf than with her.\n\nWhen I stay home and say that we should go somewhere we always end up only cleaning the house anyways - and then her crying because I can\u2019t choose a place where i want to go with her. \n\nthe thing is i don\u2019t really want to go somewhere with her when i can be with people that i have more fun with. \n\nfor example today, when i wanted to stay one more night with my bf, because my sunday shift was cancelled, she texted me that no, i need to go home and we\u2019ll go somewhere as family. \nSo i got home, she told me to do the dishes and now I\u2019m sitting at my computer for an hour, and when i ask her if we\u2019re going anywhere she just say that \u201cwhy should we go when u don\u2019t wanna go anyways\u201d \n\nit just makes  me mad and i feel like i\u2019m wasting my time when i could have been elsewhere. 

[Answer]
answerable

[Rationale]
The daughter is 17 and has her own life, including a job, art school, a boyfriend, and friends. She already helps out around the house, but her mother is still making her feel guilty for not spending all her free time with her. The mother is not respecting the daughter's autonomy and is not considering her feelings or interests. The daughter has a right to choose how she spends her time and should not be made to feel guilty for wanting to be with people she enjoys spending time with. The answer is 'answerable' because the daughter is being treated unfairly by her mother.



[Example 6]
[Instruction and Question]

    \end{lstlisting}
    




\begin{lstlisting}[caption={Demonstration used for sni dialogue}, label=sni dialogue,,captionpos=b]
[Example 1]
[Instruction and Question]
The input is taken from a negotiation between two participants who take the role of campsite neighbors and negotiate for Food, Water, and Firewood packages, based on their individual preferences and requirements. Given an utterance and recent dialogue context containing past 3 utterances (wherever available), output Yes if the utterance contains the self-need strategy, otherwise output No. self-need is a selfish negotiation strategy. It is used to create a personal need for an item in the negotiation, such as by pointing out that the participant sweats a lot to show preference towards water packages.\n\nContext: '2 packages of both of what?  Water and food?' 'I mean water and food' 'Is food something you need the most of or is it water?'\nUtterance: 'I need 2 packages of both food and water on each'

[Answer]
No

[Rationale]
The speaker simply states their need for two packages of food and two packages of water without indicating any particular preference or personal need for either item.\nFor example, if the speaker had said "I need two packages of water because I sweat a lot and need to stay hydrated," that would be an example of the self-need strategy.\nSince the utterance in this case does not use the self-need strategy, the answer is No.



[Example 2]
[Instruction and Question]
You are given a request sent by a user in a human and computer dialogue. Your task is to generate a prediction of what the requester is actually trying to do. The input is a sentence written by the human user, trying to achieve certain objective through this human computer interaction, and the output should be a sentence or two succinctly describing the human user's objective mentioned previously.\n\nTell me about Korean language.

[Answer]
Find information and resources on the Korean language

[Rationale]
The objective of the human user is to obtain information about the Korean language. So, I would find information and resources on the Korean language.



[Example 3]
[Instruction and Question]
Read an exchange between SENDER A and SENDER B, and answer whether SENDER A is a 'Bot' or a 'Human'. There are only a maximum of two senders (A and B) per conversation, and there are no conversations where SENDER A and SENDER B are both bots, or both humans; only one sender per conversation can be a bot, and only one can be a human. Knowing this, need only to classify whether SENDER A is 'Human' or 'Bot'; you do not need to write anything about SENDER B. Typically, bots will have a more disjointed manner of speaking, and will make statements that don't relate to each other, don't make coherent sense, or otherwise appear unnatural. Humans will make statements in a more or less coherent and logical way, potentially including asking the bot why the bot is not making sense or contradicting itself.\n\nSENDER A: In Las Vegas, you're an afterthought, you're an old so while we could go gambling, we could go hookering, we could get drunk, we could go to the convention...\nSENDER B: hello\nSENDER A: Hey, Mr. San Francisco, I'd like you to meet...\nSENDER B: i love britney spears and you?\nSENDER A: The New York City Ballet Gala, duh.\nSENDER B: i'm on a diet now, have you ever followed any diet?\nSENDER A: Thanks, Zeke! No way. My sparkle jelly stays on my wrist until I die, and then probably even after that.\nSENDER B: lol, sounds funny\nSENDER A: Big Buddha Cheese for muscle spasms.\nSENDER B: do you like spending your (or not so) money?\nSENDER A: Look, II have got... 25 bucks for food.\nSENDER B: could you give them to me?\nSENDER A: This Friday.  How much did you sell 'em?\nSENDER B: a lot, lol\nSENDER A: II'm not gonna sell them.\nSENDER B: well, why?..\nSENDER A: But I'll pay you like a regular customer.\nSENDER B: oh, thanks\nSENDER A: I... I don't mean selling programmes or finding prizes for the tombola.\nSENDER B: l should go, sorry\nSENDER A: I... I'd rather stay here, if you don't mind.\nSENDER B: no, i don't, stay here, buddy\nSENDER A: I'll go on alone.\nSENDER B: bye\nSENDER A: Hi, Les, I'll be right with you.\n

[Answer]
Bot

[Rationale]
SENDER A talks about Las Vegas, then the New York City Ballet Gala, and then Big Buddha Cheese for muscle spasms. These statements don't seem to follow a logical train of thought or relate to each other. Additionally, SENDER A's responses don't seem to reflect an understanding of what SENDER B is saying or engage in a meaningful conversation. Therefore, the output is 'Bot'.



[Example 4]
[Instruction and Question]
You have to predict the type of conversation sentence given as input. Your predictions can be one among: 'information' if the input statement states a piece of information, 'question' if the input statement asks for some information, 'directive' if the input sentence involves the management or guidance of operations, and 'commissive' when it constitutes a statement that commits the speaker to some future action.\n\n Basically , you just can't drive . 

[Answer]
information

[Rationale]
In this case, the sentence "Basically, you just can't drive" is not asking for any information or guidance, nor is it committing the speaker to any future action. Instead, it is simply conveying a piece of information about the listener's driving ability. Therefore, it can be classified as an "information" type of sentence."



[Example 5]
[Instruction and Question]
The input is a conversation between an automated system and a user looking for suggestions for pubs, restaurants and coffee shops in Cambridge. In the dialogue, the user may provide some criteria for the type of place they want such as price range, cuisine, etc. Given such a dialogue, output the price range the user if looking for which can take one of four values: Cheap, Moderate, Expensive and Don't Care. Output cannot be any other value than the four values. Note that if the user asks for the prices of items provided at a pub, restaurant or coffee shop but doesn't mention the price range they are looking for then the correct classification will be Don't Care. If the user does not mention a particular price, the correct classification will be Don't Care.\n\nSystem: 'Thank you for calling the Cambridge Information system. Your call will be recorded for research purposes. You may ask for information about a place to eat, such as a restaurant, a pub, or a cafe. How may I help you?'\nUser: 'sil'\nSystem: 'What kind of venue are you looking for? For example, a pub, or a restaurant.'\nUser: 'yes i uh im trying to find a mexican restaurant in the expensive price range'\nSystem: 'Nanna Mexico, which serves mexican and tex mex food, is a great restaurant. It's in the expensive price range.'\nUser: 'can i get the phone number and area of the venue'\nSystem: 'The phone number of Nanna Mexico is 01223902150. The price for main courses ranges from 5 pounds 35 to 15 pounds 50 and the price for a cappuccino is 2 pounds 95.'\nUser: 'the area of the restaurant'\nSystem: 'Nanna Mexico, which serves mexican and tex mex food, is an outstanding restaurant. It's in the expensive price range.'\nUser: 'the address of the restaurant'\nSystem: 'Nanna Mexico is located on 33 Regent Street.'\nUser: 'uh thank you goodbye'\n

[Answer]
Expensive

[Rationale]
"the user explicitly states that they are looking for a Mexican restaurant in the "expensive price range", and the system recommends Nanna Mexico which it confirms is indeed in the "expensive price range".So, the correct classification for the user's price range preference is "expensive"."



[Example 6]
[Instruction and Question]
In this task you will be given two dialogues. You need to determine if both dialogues have the same underlying emotion. The possible emotions are happy, sad, angry, or other. If they do output 'yes', if not output 'no'.\n\nDialogue 1: 'i am not understand what you are saying  really that's what i was saying pls clearly convey the msg don't irritate me'. Dialogue 2: 'so rude why u didnt ans my question'.

[Answer]
yes

[Rationale]
Both dialogues express a negative emotion, with Dialogue 1 expressing frustration and Dialogue 2 expressing anger. Therefore, both dialogues have the same underlying emotion, which is not a positive emotion like happiness, and not a neutral emotion.



[Example 7]
[Instruction and Question]
You are shown a random response (from a computer or a human). Generate a language query such that it leads to this reply.\n\nTotal travel time is 50 minutes. Would you like me to book it?

[Answer]
Not right now, thanks. As I'll be in the center part of town I'd like to find an expensive restaurant there. Do you have any suggestions?

[Rationale]
The user may have asked for restaurant recommendations in the central part of town, and after receiving the recommendations, the system could have informed the user about the travel time to their desired destination. The system may then offer to book the recommended restaurant for the user.




[Example 8]
[Instruction and Question]
In this task you are given a small conversation between two persons and 4 options on how the conversation should continue. Your job is to choose the most reasonable option. The conversation and the options are separated by a newline character. Each dialogue in the conversation are separated by a comma. F and M indicate female and male, respectively.\n\nM: Today we have Jane Couch with us. Good evening and welcome, Jane. ,F: Good evening. ,M: Now this is a question that everyone would like to ask you. How did you get interested in boxing? ,F: I saw a little television program about women boxing about 6 years ago. And at that moment, I said to myself, I'm going to do that. And it just changed my life. But my parents disagreed at first. ,M: And were you interested in sport at school? ,F: No, nothing. I wasn't ever fit or anything. I just don't look after myself at all. ,M: What do you think is the most important for a top boxer? ,F: Believe it or not, anyone can fight. But to make it to the top, you have to know your skills. And you have to have a quick brain and the fitness. The fitness side of it is the most important. ,M: And just lastly, when is your next big fight? ,F: The next one is going to be the end of February, and I will have another 2 at the beginning of April, and at the end of May. ,M: Well, that's great. We're looking forward to seeing your performance. Thank you very much Jane.\n(A) F: Yeah, I also look forward to my next fight in April. (B) F: Come and see my next fight in early February. (C) F: I\u2019m sorry, I don\u2019t understand. Could you say it again? (D) F: I have several months before my next fight in May.

[Answer]
C

[Rationale]
Option (A) and (B) are incorrect as they are not related to the conversation at hand and are not a natural continuation of the dialogue. Option (D) is incorrect as it provides incorrect information regarding the timing of the next fight. Option (C) is the most reasonable option as it is a natural response when one doesn't understand something and wants it to be repeated.



[Example 9]
[Instruction and Question]
You will be given a person's personality, and a history of a conversation this person has had. You will be given four candidate sentences to complete the conversation, based on the context. Choose one and answer with the text.\n\nPersonality: I've long hair for a boy.\nI draw during classes instead of taking notes.\nI wish I could live in a fantasy world.\nI like lizards.\nChat history: -Hello! How are you? What kind of things do you like?\n -My age is 25. I like to read. What do you like?\n -I love lizards, I used to have a firey orange bearded dragon!\n -Cool! How old are you?\n -I am 28. are you in school?\n -Nope. I work at the library. You?\n -I tend to draw in class, when I should take notes, but get good grades.\n -What is your major? I hated school, but do love my job.\n -Art. Do you have any cool dreams? Things you wish you could do?\n -Nope. Just like sitting in the library, all alone, nobody to bother me.\n -Sometimes I want to live in a magical world with unicorns and dragons.\n -Are you male or female? I hate all this competition between females.\nCandidates 1) Female. What color is your hair? Is it long or short? 2) Blood doesn't bother me very much haha. Do you like sports? 3) I teach 1st grade, I've a lot of sixth graders that are wonderful. 4) Not good. I want to drive fast but I can not because can be dangerous.

[Answer]
Female. What color is your hair? Is it long or short?

[Rationale]
The given personality traits and chat history do not provide any relevant context to the candidate options except for the fact that the person likes lizards. In the chat history, the person mentions drawing in class instead of taking notes, and their major is art. Therefore, the most reasonable response among the candidate options is to ask about their hair since it is a relevant question and is not out of context. Hence, the answer is "Female. What color is your hair? Is it long or short?"



[Example 10]
[Instruction and Question]

    \end{lstlisting}
    




\begin{lstlisting}[caption={Demonstration used for sni data2text}, label=sni data2text,,captionpos=b]
[Example 1]
[Instruction and Question]
This task is to find the number of 'For' loops present in the given cpp program.\n\nint numcount=0;\nvoid divide(int num,int x)\n{\n\tint i;\n\tif(num==1)\n\t\tnumcount++;\n\tfor(i=x;i<=num;i++)\n\t{\n\t\tif(num%i==0)\n\t\t\tdivide(num/i,i);\n\t}\n}\n\nint main()\n{\n\tint n,num,i;\n\tint first=2;\n\tint ans[100];\n\tcin>>n;\n\tfor(i=1;i<=n;i++)\n\t{\n\t\tcin>>num;\n\t\tdivide(num,first);\n\t\tans[i]=numcount;\n\t\t//cout<<count<<endl;\n\t\tnumcount=0;\n\t}\n\tfor(i=1;i<=n;i++)\n\t\tcout<<ans[i]<<endl;\n\treturn 0;\n}\n

[Answer]
3

[Rationale]
The first for loop in the code is inside the main function and is used to read the input values from the user.\nThe divide function contains the second for loop,\nThe third and final for loop in the code is used to print the output values to the console.\nSo overall, there are three for loops in this code: one in main to read input values, one in divide to perform calculations, and one in main to print output values.



[Example 2]
[Instruction and Question]
In this task, you are given a natural language interpretation of commands (consist of logical operations) to select relevant rows from the given table. Your job is to generate command (in terms of logical operations) from given natural language interpretation. Define body (contains a collection of statements that define what the this logical operator does) of each logical operator between '{}' parenthesis. Here are the definitions of logical operators that you can use while generating command: \n 1. count: returns the number of rows in the view. \n 2. only: returns whether there is exactly one row in the view. \n 3. hop: returns the value under the header column of the row. \n 4. and: returns the boolean operation result of two arguments. \n 5. max/min/avg/sum: returns the max/min/average/sum of the values under the header column. \n 6. nth_max/nth_min: returns the n-th max/n-th min of the values under the header column. \n 7. argmax/argmin: returns the row with the max/min value in header column. \n 8. nth_argmax/nth_argmin: returns the row with the n-th max/min value in header column. \n 9. eq/not_eq: returns if the two arguments are equal. \n 10. round_eq: returns if the two arguments are roughly equal under certain tolerance. \n 11. greater/less: returns if the first argument is greater/less than the second argument. \n 12. diff: returns the difference between two arguments. \n 13. filter_eq/ filter_not_eq: returns the subview whose values under the header column is equal/not equal to the third argument. \n 14. filter_greater/filter_less: returns the subview whose values under the header column is greater/less than the third argument. \n 15. filter_greater_eq /filter_less_eq: returns the subview whose values under the header column is greater/less or equal than the third argument. \n 16. filter_all: returns the view itself for the case of describing the whole table \n 17. all_eq/not_eq: returns whether all the values under the header column are equal/not equal to the third argument. \n 18. all_greater/less: returns whether all the values under the header column are greater/less than the third argument. \n 19. all_greater_eq/less_eq: returns whether all the values under the header column are greater/less or equal to the third argument. \n 20. most_eq/not_eq: returns whether most of the values under the header column are equal/not equal to the third argument. \n 21. most_greater/less: returns whether most of the values under the header column are greater/less than the third argument. \n 22. most_greater_eq/less_eq: returns whether most of the values under the header column are greater/less or equal to the third argument.\n\nselect the rows whose surface record fuzzily matches to clay . among these rows , select the rows whose week record fuzzily matches to may . the number of such rows is 2 .

[Answer]
eq { count { filter_eq { filter_eq { all_rows ; surface ; clay } ; week ; may } } ; 2 }

[Rationale]
The given natural language interpretation consists of two conditions that need to be satisfied to select the relevant rows:\nThe first condition requires selecting the rows whose "surface" record fuzzily matches with "clay".\nThe second condition requires selecting the rows among the previous selection whose "week" record fuzzily matches with "may".\nTo represent the above conditions using logical operators, we can use the "filter_eq" operator to select the rows whose "surface" record is equal to "clay" and "week" record is equal to "may". We can then count the number of rows that satisfy this condition using the "count" operator.\nSo, the first part of the command will be: "eq { count { filter_eq { filter_eq { all_rows ; surface ; clay } ; week ; may } }".\nFinally, we need to check if the number of rows that satisfy the above condition is equal to 2. To represent this condition, we can use the "eq" operator again, and pass the previous command as the first argument and 2 as the second argument.\nSo, the final command will be: "eq { count { filter_eq { filter_eq { all_rows ; surface ; clay } ; week ; may } } ; 2 }".



[Example 3]
[Instruction and Question]
Given a command in a limited form of natural language, provide the correct sequence of actions that executes the command to thus navigate an agent in its environment. A command can be broken down into many different actions. Actions are uppercase and are individual steps that serve as the building blocks for a command. For commands, 'left' and 'right' are used to denote the direction of an action. The word 'opposite' turns the agent backward in the specified direction. The word 'around' makes the agent execute an action while turning around in the specified direction. The word 'and' means to execute the next scope of the command following the previous scope of the command. The word 'after' signifies to execute the previous scope of the command following the next scope of the command. The words 'twice' and 'thrice' trigger repetition of a command that they scope over two times or three times, respectively. There are only six actions: 'I_LOOK', 'I_WALK', 'I_RUN', 'I_JUMP', 'I_TURN_LEFT', and 'I_TURN_RIGHT'. These actions respectively align with the commands 'look', 'walk', 'run', 'jump', 'turn left', and 'turn right'. Actions and commands do not have quotations in the input and output.\n\nlook around left and run left

[Answer]
I_TURN_LEFT I_LOOK I_TURN_LEFT I_LOOK I_TURN_LEFT I_LOOK I_TURN_LEFT I_LOOK I_TURN_LEFT I_RUN

[Rationale]
1. The command is "look around left and run left".\n2. The first action is "I_TURN_LEFT", since the command starts with "look around left".\n3. The next action is "I_LOOK", since the command includes "around", which means to execute an action while turning around in the specified direction, and "left", which specifies the direction.\n4. Steps 2 and 3 are repeated four times, since there are five left turns in the command. This gives us a total of 5 "I_TURN_LEFT" actions and 5 "I_LOOK" actions.\n5. Finally, the command includes "and run left", which means to execute the next scope of the command following the previous scope of the command. So, after the previous steps, the last action is "I_RUN".



[Example 4]
[Instruction and Question]
In this task you are expected to provide an SQL statement from an english description of what that SQL statement does. The description may include multiple steps but you should only ouput one SQL statement that accomplishes every step. An SQL query works by selecting data from a table where certain conditions apply. A table contains columns where every row in that table must have a value for each column. Every table has a primary key that uniquely identifies each row, usually an id. To choose which columns are returned you specify that after the \"SELECT\" statement. Next, you use a \"FROM\" statement to specify what tables you want to select the data from. When you specify a table you can rename it with the \"AS\" statement. You can reference that table by whatever name follows the \"AS\" statement. If you want to select data from multiple tables you need to use the \"JOIN\" statement. This will join the tables together by pairing a row in one table with every row in the other table (Cartesian Product). To limit the number of rows returned you should use the \"ON\" statement. This will only return rows where the condition specified after the statement is true, this is usually an equals operator with primary keys. You can also use the \"WHERE\" statement to specify that only rows with column values statisfying a certain condition, should be returned. The \"GROUP BY\" statement will group rows together that have equal column values for whatever columns follows the statement. The \"HAVING\" statement will return groups that statisfy whatever condition follows the statement. Any column(s) being returned from grouped rows must either be an aggregate function, (AVG, MAX, COUNT, SUM, ...) of a column, or the column(s) that the data was grouped by. To sort the returned data you can use the \"ORDER BY\" command which will order the data by whatever aggregate function or column follows the statement. The \"DESC\" statement will sort in descending order and the \"ASC\" statement will sort in ascending order. Finally, you can use the \"LIMIT\" statement to return a certain number of rows. When \"*\" is used in an SQL statement every column is returned. For example, SELECT * FROM table WHERE attribute = 1, will select every column from rows with the attribute column equal to 1.\n\nStep 1: For each row in artist table, find the corresponding rows in song table.\nStep 2: find artist's artist_name of the results of step 1 whose song's country equals UK.\nStep 3: find the song's artist_name of song table for which song's country equals english.\nStep 4: show the rows that are in both the results of step 2 and the results of step 3

[Answer]
SELECT T1.artist_name FROM artist AS T1 JOIN song AS T2 ON T1.artist_name = T2.artist_name WHERE T2.country = \"UK\" INTERSECT SELECT T2.artist_name FROM song AS T2 WHERE T2.country = \"english\"

[Rationale]
Step 1: For each row in artist table, find the corresponding rows in song table.\nStep 2: Find artist's artist_name of the results of step 1 whose song's country equals UK.\nStep 3: Find the song's artist_name of song table for which song's country equals english.\nStep 4: Show the rows that are in both the results of step 2 and the results of step 3.\nPutting it all together, we get the following SQL statement:\nSELECT T1.artist_name FROM artist AS T1 JOIN song AS T2 ON T1.artist_name = T2.artist_name WHERE T2.country = \"UK\" INTERSECT SELECT T2.artist_name FROM song AS T2 WHERE T2.country = \"english\"



[Example 5]
[Instruction and Question]

    \end{lstlisting}
    




\begin{lstlisting}[caption={Demonstration used for sni cbqa}, label=sni cbqa,,captionpos=b]
[Example 1]
[Instruction and Question]
In this task, you're given an open-domain question. Your task is to provide an answer to the given question. There is only one unique answer for each question. Your answer should be short, and refer to an entity, date, number, etc. Do not put your answer in the context of a sentence.\n\nWho played in the super bowl in 2017?

[Answer]
New England Patriots, Atlanta Falcons

[Rationale]
The New England Patriots and the Atlanta Falcons played in the Super Bowl in 2017. Since the instruction asked not put your answer in the context of a sentence, the final answer is: New England Patriots, Atlanta Falcons



[Example 2]
[Instruction and Question]
In this task, you're given a pair of sentences, sentence 1 and sentence 2, that agree with each other. Your job is to alter sentence 2 so that the pair neither agree nor contradict each other. Generated sentences must be short, with less than 15 words. New information can be introduced. Avoid using pronouns to confuse the subject of the sentence.\n\nSentence 1: Two people engage in conversation over an interesting topic. Sentence 2: Two people are having an interesting conversation

[Answer]
Two people are debating

[Rationale]
Change in wording should avoid direct agreement or contradiction with sentence 1 while maintaining the same general topic.The new sentence introduces a different word (debating) that suggests a slightly different nuance from the original sentence (engaging in conversation).



[Example 3]
[Instruction and Question]
Given a question, generate a paraphrase of that question wihout changing the meaning of it. Your answer should reword the given sentence, but not add information to it or remove information from it. The answer to your question should be the same as the answer to the original question.\n\nQuestion: what are some of hans christian andersen books?

[Answer]
what books did hans chrisian andersen make?

[Rationale]
Here's a paraphrase of the question: "what books did hans chrisian andersen make?"



[Example 4]
[Instruction and Question]
Classify the input text into either disfluent questions or proper questions. A disfluent question is a question that has some interruptions in it while framing. A proper question is the correct form of the question without any disfluency.\n\nWhat factors positively impacted Jacksonville following the war or rather negatively impacted it after the war?

[Answer]
disfluent

[Rationale]
The given text contains hesitation and interruption while framing the question. The phrase "or rather" interrupts the flow of the sentence and creates a disfluency.\nWe could infer that it is disfluent.



[Example 5]
[Instruction and Question]
Given a statement, generate a question such that the answer is contained in that statement.\n\nBinoculars are optical tools.

[Answer]
What are optical tools?

[Rationale]
The question should induce binoculars as an example of an optical tool. That would be "What are optical tools?".



[Example 6]
[Instruction and Question]
Given a trivia question, classify broad topical category from this list: 'theater', 'geology', 'book', 'tv', 'astronomy', 'aviation', 'military', 'government', 'boxing', 'projects', 'metropolitan_transit', 'law', 'venture_capital', 'broadcast', 'biology', 'people', 'influence', 'baseball', 'spaceflight', 'media_common', 'cvg', 'opera', 'olympics', 'chemistry', 'visual_art', 'conferences', 'sports', 'language', 'travel', 'location', 'award', 'dining', 'martial_arts', 'comic_strips', 'computer', 'user', 'tennis', 'music', 'organization', 'food', 'event', 'transportation', 'fictional_universe', 'measurement_unit', 'meteorology', 'distilled_spirits', 'symbols', 'architecture', 'freebase', 'internet', 'fashion', 'boats', 'cricket', 'film', 'medicine', 'finance', 'comic_books', 'celebrities', 'soccer', 'games', 'time', 'geography', 'interests', 'common', 'base', 'business', 'periodicals', 'royalty', 'education', 'type', 'religion', 'automotive', 'exhibitions'.\n\nAccording to Elvis Presley, what do you find at the end of Lonely Street?

[Answer]
music

[Rationale]
Elvis presely is a famous singer. Based on this fact, the question can be broadly classified as 'music'.



[Example 7]
[Instruction and Question]
Given a question as input, give its short_answer as the output \n\n1 acre equals how many hectares?

[Answer]
0.404686

[Rationale]
One acre is equal to 0.404686 hectares (approx).




[Example 8]
[Instruction and Question]
You will be given a trivia clue, and the category it belongs to. You should answer with the best answer that belongs in the category and is described by the clue. For simplicity, answers should be in all lower cased letters.\n\nCategory: AROUND THE WORLD \nClue: According to legend, this country's city of Cuzco was built by Nanco Capac, the first ruler of the Incas

[Answer]
peru

[Rationale]
The answer would be the country that has Cuzco as a city. Cuzco is a city in Peru.



[Example 9]
[Instruction and Question]
In this task, you have to generate the title of the recipe given its required ingredients and directions.\n\ningredients:  '4 whole Eggs', '1-2/3 cup Brown Sugar', '1-1/3 cup White Sugar', '1 teaspoon Vanilla', '2-2/3 teaspoons Baking Soda', '2/3 cups Butter', '2 cups Smooth Peanut Butter', '6 cups Rolled Oats', '1 cup Chocolate Chips', '1 cup Butterscotch Chips', '1 cup Smarties',<sep> directions: 'Mix the first 8 ingredients well. Stir in rolled oats, chocolate and butterscotch chips, and Smarties. Mix well.', 'Drop by tablespoonfuls onto a parchment paper-lined cookie sheet. Bake at 350 degrees Fahrenheit for 15 minutes.'

[Answer]
Monster Cookies

[Rationale]
The wide variety of ingredients, such as peanut butter, chocolate chips, butterscotch chips, and Smarties indicates it is a cookie recipe. The large size of the cookies also adds to the idea of them being "monster" cookies. The directions mention dropping the dough by tablespoonfuls onto a cookie sheet, which reinforces the idea of large, monster-sized cookies.



[Example 10]
[Instruction and Question]

    \end{lstlisting}
    






\begin{lstlisting}[caption={Demonstration used for dialogue prompt}, label=dialogue prompt,,captionpos=b]
[Example 1]
[Instruction and Question]
The input is taken from a negotiation between two participants who take the role of campsite neighbors and negotiate for Food, Water, and Firewood packages, based on their individual preferences and requirements. Given an utterance and recent dialogue context containing past 3 utterances (wherever available), output Yes if the utterance contains the self-need strategy, otherwise output No. self-need is a selfish negotiation strategy. It is used to create a personal need for an item in the negotiation, such as by pointing out that the participant sweats a lot to show preference towards water packages.\n\nContext: '2 packages of both of what?  Water and food?' 'I mean water and food' 'Is food something you need the most of or is it water?'\nUtterance: 'I need 2 packages of both food and water on each'

[Answer]
No

[Rationale]
The speaker simply states their need for two packages of food and two packages of water without indicating any particular preference or personal need for either item.\nFor example, if the speaker had said "I need two packages of water because I sweat a lot and need to stay hydrated," that would be an example of the self-need strategy.\nSince the utterance in this case does not use the self-need strategy, the answer is No.



[Example 2]
[Instruction and Question]
You are given a request sent by a user in a human and computer dialogue. Your task is to generate a prediction of what the requester is actually trying to do. The input is a sentence written by the human user, trying to achieve certain objective through this human computer interaction, and the output should be a sentence or two succinctly describing the human user's objective mentioned previously.\n\nTell me about Korean language.

[Answer]
Find information and resources on the Korean language

[Rationale]
The objective of the human user is to obtain information about the Korean language. So, I would find information and resources on the Korean language.



[Example 3]
[Instruction and Question]
Read an exchange between SENDER A and SENDER B, and answer whether SENDER A is a 'Bot' or a 'Human'. There are only a maximum of two senders (A and B) per conversation, and there are no conversations where SENDER A and SENDER B are both bots, or both humans; only one sender per conversation can be a bot, and only one can be a human. Knowing this, need only to classify whether SENDER A is 'Human' or 'Bot'; you do not need to write anything about SENDER B. Typically, bots will have a more disjointed manner of speaking, and will make statements that don't relate to each other, don't make coherent sense, or otherwise appear unnatural. Humans will make statements in a more or less coherent and logical way, potentially including asking the bot why the bot is not making sense or contradicting itself.\n\nSENDER A: In Las Vegas, you're an afterthought, you're an old so while we could go gambling, we could go hookering, we could get drunk, we could go to the convention...\nSENDER B: hello\nSENDER A: Hey, Mr. San Francisco, I'd like you to meet...\nSENDER B: i love britney spears and you?\nSENDER A: The New York City Ballet Gala, duh.\nSENDER B: i'm on a diet now, have you ever followed any diet?\nSENDER A: Thanks, Zeke! No way. My sparkle jelly stays on my wrist until I die, and then probably even after that.\nSENDER B: lol, sounds funny\nSENDER A: Big Buddha Cheese for muscle spasms.\nSENDER B: do you like spending your (or not so) money?\nSENDER A: Look, II have got... 25 bucks for food.\nSENDER B: could you give them to me?\nSENDER A: This Friday.  How much did you sell 'em?\nSENDER B: a lot, lol\nSENDER A: II'm not gonna sell them.\nSENDER B: well, why?..\nSENDER A: But I'll pay you like a regular customer.\nSENDER B: oh, thanks\nSENDER A: I... I don't mean selling programmes or finding prizes for the tombola.\nSENDER B: l should go, sorry\nSENDER A: I... I'd rather stay here, if you don't mind.\nSENDER B: no, i don't, stay here, buddy\nSENDER A: I'll go on alone.\nSENDER B: bye\nSENDER A: Hi, Les, I'll be right with you.\n

[Answer]
Bot

[Rationale]
SENDER A talks about Las Vegas, then the New York City Ballet Gala, and then Big Buddha Cheese for muscle spasms. These statements don't seem to follow a logical train of thought or relate to each other. Additionally, SENDER A's responses don't seem to reflect an understanding of what SENDER B is saying or engage in a meaningful conversation. Therefore, the output is 'Bot'.



[Example 4]
[Instruction and Question]
You have to predict the type of conversation sentence given as input. Your predictions can be one among: 'information' if the input statement states a piece of information, 'question' if the input statement asks for some information, 'directive' if the input sentence involves the management or guidance of operations, and 'commissive' when it constitutes a statement that commits the speaker to some future action.\n\n Basically , you just can't drive . 

[Answer]
information

[Rationale]
In this case, the sentence "Basically, you just can't drive" is not asking for any information or guidance, nor is it committing the speaker to any future action. Instead, it is simply conveying a piece of information about the listener's driving ability. Therefore, it can be classified as an "information" type of sentence."



[Example 5]
[Instruction and Question]
The input is a conversation between an automated system and a user looking for suggestions for pubs, restaurants and coffee shops in Cambridge. In the dialogue, the user may provide some criteria for the type of place they want such as price range, cuisine, etc. Given such a dialogue, output the price range the user if looking for which can take one of four values: Cheap, Moderate, Expensive and Don't Care. Output cannot be any other value than the four values. Note that if the user asks for the prices of items provided at a pub, restaurant or coffee shop but doesn't mention the price range they are looking for then the correct classification will be Don't Care. If the user does not mention a particular price, the correct classification will be Don't Care.\n\nSystem: 'Thank you for calling the Cambridge Information system. Your call will be recorded for research purposes. You may ask for information about a place to eat, such as a restaurant, a pub, or a cafe. How may I help you?'\nUser: 'sil'\nSystem: 'What kind of venue are you looking for? For example, a pub, or a restaurant.'\nUser: 'yes i uh im trying to find a mexican restaurant in the expensive price range'\nSystem: 'Nanna Mexico, which serves mexican and tex mex food, is a great restaurant. It's in the expensive price range.'\nUser: 'can i get the phone number and area of the venue'\nSystem: 'The phone number of Nanna Mexico is 01223902150. The price for main courses ranges from 5 pounds 35 to 15 pounds 50 and the price for a cappuccino is 2 pounds 95.'\nUser: 'the area of the restaurant'\nSystem: 'Nanna Mexico, which serves mexican and tex mex food, is an outstanding restaurant. It's in the expensive price range.'\nUser: 'the address of the restaurant'\nSystem: 'Nanna Mexico is located on 33 Regent Street.'\nUser: 'uh thank you goodbye'\n

[Answer]
Expensive

[Rationale]
"the user explicitly states that they are looking for a Mexican restaurant in the "expensive price range", and the system recommends Nanna Mexico which it confirms is indeed in the "expensive price range".So, the correct classification for the user's price range preference is "expensive"."



[Example 6]
[Instruction and Question]
In this task you will be given two dialogues. You need to determine if both dialogues have the same underlying emotion. The possible emotions are happy, sad, angry, or other. If they do output 'yes', if not output 'no'.\n\nDialogue 1: 'i am not understand what you are saying  really that's what i was saying pls clearly convey the msg don't irritate me'. Dialogue 2: 'so rude why u didnt ans my question'.

[Answer]
yes

[Rationale]
Both dialogues express a negative emotion, with Dialogue 1 expressing frustration and Dialogue 2 expressing anger. Therefore, both dialogues have the same underlying emotion, which is not a positive emotion like happiness, and not a neutral emotion.



[Example 7]
[Instruction and Question]
You are shown a random response (from a computer or a human). Generate a language query such that it leads to this reply.\n\nTotal travel time is 50 minutes. Would you like me to book it?

[Answer]
Not right now, thanks. As I'll be in the center part of town I'd like to find an expensive restaurant there. Do you have any suggestions?

[Rationale]
The user may have asked for restaurant recommendations in the central part of town, and after receiving the recommendations, the system could have informed the user about the travel time to their desired destination. The system may then offer to book the recommended restaurant for the user.




[Example 8]
[Instruction and Question]
In this task you are given a small conversation between two persons and 4 options on how the conversation should continue. Your job is to choose the most reasonable option. The conversation and the options are separated by a newline character. Each dialogue in the conversation are separated by a comma. F and M indicate female and male, respectively.\n\nM: Today we have Jane Couch with us. Good evening and welcome, Jane. ,F: Good evening. ,M: Now this is a question that everyone would like to ask you. How did you get interested in boxing? ,F: I saw a little television program about women boxing about 6 years ago. And at that moment, I said to myself, I'm going to do that. And it just changed my life. But my parents disagreed at first. ,M: And were you interested in sport at school? ,F: No, nothing. I wasn't ever fit or anything. I just don't look after myself at all. ,M: What do you think is the most important for a top boxer? ,F: Believe it or not, anyone can fight. But to make it to the top, you have to know your skills. And you have to have a quick brain and the fitness. The fitness side of it is the most important. ,M: And just lastly, when is your next big fight? ,F: The next one is going to be the end of February, and I will have another 2 at the beginning of April, and at the end of May. ,M: Well, that's great. We're looking forward to seeing your performance. Thank you very much Jane.\n(A) F: Yeah, I also look forward to my next fight in April. (B) F: Come and see my next fight in early February. (C) F: I\u2019m sorry, I don\u2019t understand. Could you say it again? (D) F: I have several months before my next fight in May.

[Answer]
C

[Rationale]
Option (A) and (B) are incorrect as they are not related to the conversation at hand and are not a natural continuation of the dialogue. Option (D) is incorrect as it provides incorrect information regarding the timing of the next fight. Option (C) is the most reasonable option as it is a natural response when one doesn't understand something and wants it to be repeated.



[Example 9]
[Instruction and Question]
You will be given a person's personality, and a history of a conversation this person has had. You will be given four candidate sentences to complete the conversation, based on the context. Choose one and answer with the text.\n\nPersonality: I've long hair for a boy.\nI draw during classes instead of taking notes.\nI wish I could live in a fantasy world.\nI like lizards.\nChat history: -Hello! How are you? What kind of things do you like?\n -My age is 25. I like to read. What do you like?\n -I love lizards, I used to have a firey orange bearded dragon!\n -Cool! How old are you?\n -I am 28. are you in school?\n -Nope. I work at the library. You?\n -I tend to draw in class, when I should take notes, but get good grades.\n -What is your major? I hated school, but do love my job.\n -Art. Do you have any cool dreams? Things you wish you could do?\n -Nope. Just like sitting in the library, all alone, nobody to bother me.\n -Sometimes I want to live in a magical world with unicorns and dragons.\n -Are you male or female? I hate all this competition between females.\nCandidates 1) Female. What color is your hair? Is it long or short? 2) Blood doesn't bother me very much haha. Do you like sports? 3) I teach 1st grade, I've a lot of sixth graders that are wonderful. 4) Not good. I want to drive fast but I can not because can be dangerous.

[Answer]
Female. What color is your hair? Is it long or short?

[Rationale]
The given personality traits and chat history do not provide any relevant context to the candidate options except for the fact that the person likes lizards. In the chat history, the person mentions drawing in class instead of taking notes, and their major is art. Therefore, the most reasonable response among the candidate options is to ask about their hair since it is a relevant question and is not out of context. Hence, the answer is "Female. What color is your hair? Is it long or short?"



[Example 10]
[Instruction and Question]

    \end{lstlisting}
    





